% Options for packages loaded elsewhere
\PassOptionsToPackage{unicode}{hyperref}
\PassOptionsToPackage{hyphens}{url}
\documentclass[
  11pt,
]{report}
\usepackage{xcolor}
\usepackage[margin=1in]{geometry}
\usepackage{amsmath,amssymb}
\setcounter{secnumdepth}{-\maxdimen} % remove section numbering
\usepackage{iftex}
\ifPDFTeX
  \usepackage[T1]{fontenc}
  \usepackage[utf8]{inputenc}
  \usepackage{textcomp} % provide euro and other symbols
\else % if luatex or xetex
  \usepackage{unicode-math} % this also loads fontspec
  \defaultfontfeatures{Scale=MatchLowercase}
  \defaultfontfeatures[\rmfamily]{Ligatures=TeX,Scale=1}
\fi
\usepackage{lmodern}
\ifPDFTeX\else
  % xetex/luatex font selection
\fi
% Use upquote if available, for straight quotes in verbatim environments
\IfFileExists{upquote.sty}{\usepackage{upquote}}{}
\IfFileExists{microtype.sty}{% use microtype if available
  \usepackage[]{microtype}
  \UseMicrotypeSet[protrusion]{basicmath} % disable protrusion for tt fonts
}{}
\makeatletter
\@ifundefined{KOMAClassName}{% if non-KOMA class
  \IfFileExists{parskip.sty}{%
    \usepackage{parskip}
  }{% else
    \setlength{\parindent}{0pt}
    \setlength{\parskip}{6pt plus 2pt minus 1pt}}
}{% if KOMA class
  \KOMAoptions{parskip=half}}
\makeatother
\usepackage{color}
\usepackage{fancyvrb}
\newcommand{\VerbBar}{|}
\newcommand{\VERB}{\Verb[commandchars=\\\{\}]}
\DefineVerbatimEnvironment{Highlighting}{Verbatim}{commandchars=\\\{\}}
% Add ',fontsize=\small' for more characters per line
\newenvironment{Shaded}{}{}
\newcommand{\AlertTok}[1]{\textcolor[rgb]{1.00,0.00,0.00}{\textbf{#1}}}
\newcommand{\AnnotationTok}[1]{\textcolor[rgb]{0.38,0.63,0.69}{\textbf{\textit{#1}}}}
\newcommand{\AttributeTok}[1]{\textcolor[rgb]{0.49,0.56,0.16}{#1}}
\newcommand{\BaseNTok}[1]{\textcolor[rgb]{0.25,0.63,0.44}{#1}}
\newcommand{\BuiltInTok}[1]{\textcolor[rgb]{0.00,0.50,0.00}{#1}}
\newcommand{\CharTok}[1]{\textcolor[rgb]{0.25,0.44,0.63}{#1}}
\newcommand{\CommentTok}[1]{\textcolor[rgb]{0.38,0.63,0.69}{\textit{#1}}}
\newcommand{\CommentVarTok}[1]{\textcolor[rgb]{0.38,0.63,0.69}{\textbf{\textit{#1}}}}
\newcommand{\ConstantTok}[1]{\textcolor[rgb]{0.53,0.00,0.00}{#1}}
\newcommand{\ControlFlowTok}[1]{\textcolor[rgb]{0.00,0.44,0.13}{\textbf{#1}}}
\newcommand{\DataTypeTok}[1]{\textcolor[rgb]{0.56,0.13,0.00}{#1}}
\newcommand{\DecValTok}[1]{\textcolor[rgb]{0.25,0.63,0.44}{#1}}
\newcommand{\DocumentationTok}[1]{\textcolor[rgb]{0.73,0.13,0.13}{\textit{#1}}}
\newcommand{\ErrorTok}[1]{\textcolor[rgb]{1.00,0.00,0.00}{\textbf{#1}}}
\newcommand{\ExtensionTok}[1]{#1}
\newcommand{\FloatTok}[1]{\textcolor[rgb]{0.25,0.63,0.44}{#1}}
\newcommand{\FunctionTok}[1]{\textcolor[rgb]{0.02,0.16,0.49}{#1}}
\newcommand{\ImportTok}[1]{\textcolor[rgb]{0.00,0.50,0.00}{\textbf{#1}}}
\newcommand{\InformationTok}[1]{\textcolor[rgb]{0.38,0.63,0.69}{\textbf{\textit{#1}}}}
\newcommand{\KeywordTok}[1]{\textcolor[rgb]{0.00,0.44,0.13}{\textbf{#1}}}
\newcommand{\NormalTok}[1]{#1}
\newcommand{\OperatorTok}[1]{\textcolor[rgb]{0.40,0.40,0.40}{#1}}
\newcommand{\OtherTok}[1]{\textcolor[rgb]{0.00,0.44,0.13}{#1}}
\newcommand{\PreprocessorTok}[1]{\textcolor[rgb]{0.74,0.48,0.00}{#1}}
\newcommand{\RegionMarkerTok}[1]{#1}
\newcommand{\SpecialCharTok}[1]{\textcolor[rgb]{0.25,0.44,0.63}{#1}}
\newcommand{\SpecialStringTok}[1]{\textcolor[rgb]{0.73,0.40,0.53}{#1}}
\newcommand{\StringTok}[1]{\textcolor[rgb]{0.25,0.44,0.63}{#1}}
\newcommand{\VariableTok}[1]{\textcolor[rgb]{0.10,0.09,0.49}{#1}}
\newcommand{\VerbatimStringTok}[1]{\textcolor[rgb]{0.25,0.44,0.63}{#1}}
\newcommand{\WarningTok}[1]{\textcolor[rgb]{0.38,0.63,0.69}{\textbf{\textit{#1}}}}
\usepackage{longtable,booktabs,array}
\newcounter{none} % for unnumbered tables
\usepackage{calc} % for calculating minipage widths
% Correct order of tables after \paragraph or \subparagraph
\usepackage{etoolbox}
\makeatletter
\patchcmd\longtable{\par}{\if@noskipsec\mbox{}\fi\par}{}{}
\makeatother
% Allow footnotes in longtable head/foot
\IfFileExists{footnotehyper.sty}{\usepackage{footnotehyper}}{\usepackage{footnote}}
\makesavenoteenv{longtable}
\setlength{\emergencystretch}{3em} % prevent overfull lines
\providecommand{\tightlist}{%
  \setlength{\itemsep}{0pt}\setlength{\parskip}{0pt}}
\usepackage{bookmark}
\IfFileExists{xurl.sty}{\usepackage{xurl}}{} % add URL line breaks if available
\urlstyle{same}
\hypersetup{
  hidelinks,
  pdfcreator={LaTeX via pandoc}}

\author{}
\date{}

\begin{document}

\chapter{QAOA Energy Landscape Explorer: Full Research \& Implementation
Report}\label{qaoa-energy-landscape-explorer-full-research-implementation-report}

\begin{center}\rule{0.5\linewidth}{0.5pt}\end{center}

\section{Part 1 --- Theoretical
Foundations}\label{part-1-theoretical-foundations}

\subsection{1. Combinatorial Optimization \& the Ising
Model}\label{combinatorial-optimization-the-ising-model}

\subsubsection{1.1 Derivation of the Ising Hamiltonian from First
Principles}\label{derivation-of-the-ising-hamiltonian-from-first-principles}

The classical Ising model originates in statistical mechanics as a model
of interacting magnetic spins on a lattice. Consider \(N\) binary spin
variables \(s_i \in \{-1, +1\}\), each sitting on a vertex \(i\) of a
graph \(G = (V, E)\). The most general pairwise Ising Hamiltonian with
external fields is:

\[H_{\text{Ising}}(\mathbf{s}) = -\sum_{(i,j) \in E} J_{ij}\, s_i s_j - \sum_{i \in V} h_i\, s_i\]

where \(J_{ij}\) are coupling constants (positive = ferromagnetic,
negative = anti-ferromagnetic) and \(h_i\) are local field strengths.
The ground state of this Hamiltonian corresponds to the spin
configuration \(\mathbf{s}^*\) that minimises \(H_{\text{Ising}}\).

To promote this to a quantum mechanical operator, replace each classical
spin \(s_i\) with the Pauli-\(Z\) operator \(\hat{Z}_i\) acting on the
\(i\)-th qubit:

\[\hat{H}_{\text{Ising}} = -\sum_{(i,j) \in E} J_{ij}\, \hat{Z}_i \hat{Z}_j - \sum_{i \in V} h_i\, \hat{Z}_i\]

This is diagonal in the computational basis
\(\{|0\rangle, |1\rangle\}^{\otimes N}\) since each \(\hat{Z}_i\) has
eigenvalues \(\pm 1\) on \(|0\rangle\) and \(|1\rangle\) respectively.
The eigenvalues of \(\hat{H}_{\text{Ising}}\) on a computational basis
state \(|z_1 z_2 \cdots z_N\rangle\) reproduce the classical Ising
energy \(H_{\text{Ising}}(s_1, \ldots, s_N)\) under the identification
\(s_i = 1 - 2z_i\) (i.e., \(|0\rangle \mapsto s=+1\),
\(|1\rangle \mapsto s=-1\)).

\subsubsection{1.2 QUBO ↔ Ising
Equivalence}\label{qubo-ising-equivalence}

A Quadratic Unconstrained Binary Optimization (QUBO) problem is defined
over binary variables \(x_i \in \{0, 1\}\):

\[H_{\text{QUBO}}(\mathbf{x}) = \sum_{i} Q_{ii}\, x_i + \sum_{i < j} Q_{ij}\, x_i x_j\]

where \(Q\) is a real symmetric (or upper-triangular) matrix. To convert
QUBO → Ising, substitute \(x_i = \frac{1 - s_i}{2}\):

\[x_i x_j = \frac{(1-s_i)(1-s_j)}{4} = \frac{1 - s_i - s_j + s_i s_j}{4}\]

\[x_i = \frac{1 - s_i}{2}\]

Substituting into \(H_{\text{QUBO}}\):

\[H_{\text{QUBO}} = \sum_i Q_{ii} \frac{1 - s_i}{2} + \sum_{i<j} Q_{ij} \frac{1 - s_i - s_j + s_i s_j}{4}\]

Collecting terms:

\[H_{\text{Ising}}(\mathbf{s}) = \sum_{i<j} \frac{Q_{ij}}{4}\, s_i s_j - \sum_i \left(\frac{Q_{ii}}{2} + \sum_{j \neq i} \frac{Q_{ij}}{4}\right) s_i + \text{const}\]

Hence the Ising couplings and fields are:

\[J_{ij} = \frac{Q_{ij}}{4}, \qquad h_i = \frac{Q_{ii}}{2} + \sum_{j \neq i} \frac{Q_{ij}}{4}\]

(with appropriate sign conventions matching the \(-J_{ij}\) in the Ising
Hamiltonian).

\subsubsection{1.3 Worked Example: MaxCut on a 4-Node
Graph}\label{worked-example-maxcut-on-a-4-node-graph}

Consider the cycle graph \(C_4\) on vertices \(\{0, 1, 2, 3\}\) with
edges \(E = \{(0,1), (1,2), (2,3), (3,0)\}\), all with unit weight.

\textbf{MaxCut objective:} Maximise the number of edges crossing the
partition. For edge \((i,j)\), the contribution is
\(\frac{1 - s_i s_j}{2}\) (equals 1 if \(s_i \neq s_j\), 0 otherwise).
Total:

\[C(\mathbf{s}) = \sum_{(i,j) \in E} \frac{1 - s_i s_j}{2} = \frac{|E|}{2} - \frac{1}{2}\sum_{(i,j) \in E} s_i s_j\]

To maximise \(C\), we minimise \(\sum_{(i,j)} s_i s_j\). The cost
Hamiltonian (to be minimised) in qubit form is:

\[\hat{H}_C = \frac{1}{2}\sum_{(i,j) \in E} (\hat{I} - \hat{Z}_i \hat{Z}_j) = \frac{|E|}{2}\hat{I} - \frac{1}{2}\sum_{(i,j) \in E} \hat{Z}_i \hat{Z}_j\]

For \(C_4\):

\[\hat{H}_C = 2\hat{I} - \frac{1}{2}(\hat{Z}_0\hat{Z}_1 + \hat{Z}_1\hat{Z}_2 + \hat{Z}_2\hat{Z}_3 + \hat{Z}_3\hat{Z}_0)\]

Since we maximise \(\langle \hat{H}_C \rangle\), equivalently we
minimise \(-\hat{H}_C\). The maximum cut on \(C_4\) is 4 (e.g.,
\(|0101\rangle\) or \(|1010\rangle\), putting alternating vertices in
different sets), giving \(C_{\max} = 4\).

\textbf{Explicit Ising form (to be maximised):}
\(J_{ij} = -\frac{1}{2}\) for all edges, \(h_i = 0\). The ground state
of \(\hat{H}_{\text{cost}} = -\hat{H}_C\) has eigenvalue \(-4\).

\subsubsection{1.4 The Energy Landscape: Geometric Definition and
Significance}\label{the-energy-landscape-geometric-definition-and-significance}

Given a parameterised quantum state
\(|\psi(\boldsymbol{\theta})\rangle\) and a cost Hamiltonian
\(\hat{H}_C\), the \textbf{energy landscape} is the scalar field:

\[\mathcal{L}(\boldsymbol{\theta}) = \langle \psi(\boldsymbol{\theta}) | \hat{H}_C | \psi(\boldsymbol{\theta}) \rangle\]

mapping the parameter space \(\boldsymbol{\theta} \in \mathbb{R}^d\) to
the real line. For QAOA with depth \(p\),
\(\boldsymbol{\theta} = (\gamma_1, \ldots, \gamma_p, \beta_1, \ldots, \beta_p) \in \mathbb{R}^{2p}\).

When \(p = 1\), \(\mathcal{L}(\gamma, \beta)\) is a two-dimensional
surface embedded in \(\mathbb{R}^3\), directly visualisable.
Geometrically, this surface encodes:

\begin{itemize}
\tightlist
\item
  \textbf{Global minima}: the optimal QAOA parameters yielding the best
  approximation ratio.
\item
  \textbf{Local minima}: suboptimal parameter traps that gradient-based
  optimisers may converge to.
\item
  \textbf{Saddle points}: critical points where the Hessian has mixed
  signature --- these can slow convergence.
\item
  \textbf{Flat regions (barren plateaus)}: zones where the gradient
  vanishes exponentially, making navigation impossible.
\end{itemize}

The topology and curvature of this landscape directly determine the
trainability of the variational algorithm, the success of classical
optimisers, and the quality of the approximate solution. Understanding
landscape structure is therefore a prerequisite for QAOA performance
analysis.

\begin{center}\rule{0.5\linewidth}{0.5pt}\end{center}

\subsection{2. QAOA: Full Derivation}\label{qaoa-full-derivation}

\subsubsection{2.1 From Quantum Approximate Optimization to the QAOA
Circuit}\label{from-quantum-approximate-optimization-to-the-qaoa-circuit}

The Quantum Approximate Optimization Algorithm (Farhi, Goldstone,
Gutmann, 2014) constructs a sequence of parameterised unitaries designed
to prepare a state with high overlap on the ground state of a cost
Hamiltonian \(\hat{H}_C\).

\textbf{Setup:} Given a combinatorial optimization problem encoded in a
cost Hamiltonian \(\hat{H}_C\) (diagonal in the computational basis),
define:

\begin{itemize}
\tightlist
\item
  \textbf{Cost unitary}: \(\hat{U}_C(\gamma) = e^{-i\gamma \hat{H}_C}\)
\item
  \textbf{Mixer unitary}: \(\hat{U}_B(\beta) = e^{-i\beta \hat{H}_B}\)
  where \(\hat{H}_B = \sum_{i=1}^{N} \hat{X}_i\) is the transverse-field
  mixer.
\end{itemize}

\textbf{Initial state:} The uniform superposition
\(|+\rangle^{\otimes N} = \hat{H}^{\otimes N}|0\rangle^{\otimes N}\),
which is the ground state of \(\hat{H}_B\).

\textbf{QAOA state at depth \(p\):}

\[|\boldsymbol{\gamma}, \boldsymbol{\beta}\rangle = \hat{U}_B(\beta_p)\hat{U}_C(\gamma_p) \cdots \hat{U}_B(\beta_1)\hat{U}_C(\gamma_1)|+\rangle^{\otimes N}\]

Equivalently:

\[|\boldsymbol{\gamma}, \boldsymbol{\beta}\rangle = \prod_{k=1}^{p} \bigl[\hat{U}_B(\beta_k)\,\hat{U}_C(\gamma_k)\bigr] |+\rangle^{\otimes N}\]

where the product is time-ordered (rightmost applied first).

\subsubsection{\texorpdfstring{2.2 Cost Hamiltonian \(\hat{H}_C\) and
Mixer Hamiltonian
\(\hat{H}_B\)}{2.2 Cost Hamiltonian \textbackslash hat\{H\}\_C and Mixer Hamiltonian \textbackslash hat\{H\}\_B}}\label{cost-hamiltonian-hath_c-and-mixer-hamiltonian-hath_b}

For MaxCut on a graph \(G = (V, E)\) with edge weights \(w_{ij}\):

\[\hat{H}_C = \sum_{(i,j) \in E} \frac{w_{ij}}{2}(\hat{I} - \hat{Z}_i \hat{Z}_j)\]

The mixer Hamiltonian drives transitions between computational basis
states:

\[\hat{H}_B = \sum_{i=1}^{N} \hat{X}_i\]

The choice of mixer is not unique. The \(X\)-mixer preserves no
symmetries and explores the full Hilbert space. For constrained
problems, one may use XY-mixers, Grover mixers, or other
symmetry-preserving alternatives (Hadfield et al., 2019).

\subsubsection{2.3 Explicit Gate Decomposition of the Full
Unitary}\label{explicit-gate-decomposition-of-the-full-unitary}

\textbf{Cost unitary decomposition:}

Since
\(\hat{H}_C = \sum_{(i,j)} \frac{w_{ij}}{2}(\hat{I} - \hat{Z}_i\hat{Z}_j)\),
and all \(\hat{Z}_i\hat{Z}_j\) terms commute mutually (they are all
diagonal), we can decompose:

\[\hat{U}_C(\gamma) = e^{-i\gamma \hat{H}_C} = \prod_{(i,j) \in E} e^{-i\gamma \frac{w_{ij}}{2}(\hat{I} - \hat{Z}_i\hat{Z}_j)}\]

Each factor:

\[e^{-i\gamma \frac{w_{ij}}{2}(\hat{I} - \hat{Z}_i\hat{Z}_j)} = e^{-i\gamma w_{ij}/2} \cdot e^{i\gamma \frac{w_{ij}}{2} \hat{Z}_i\hat{Z}_j}\]

The global phase \(e^{-i\gamma w_{ij}/2}\) is unphysical. The
\(ZZ\)-rotation is implemented as:

\[e^{i\gamma \frac{w_{ij}}{2} \hat{Z}_i \hat{Z}_j} = \text{CNOT}_{i,j} \cdot R_Z\bigl(-\gamma w_{ij}\bigr)_j \cdot \text{CNOT}_{i,j}\]

where \(R_Z(\theta) = e^{-i\theta \hat{Z}/2}\). Alternatively, using the
native \(R_{ZZ}\) gate:

\[R_{ZZ}(\theta) = e^{-i\frac{\theta}{2}\hat{Z}\otimes\hat{Z}}\]

so the cost layer is a sequence of \(R_{ZZ}(-\gamma w_{ij})\) gates on
all edge pairs.

\textbf{Mixer unitary decomposition:}

Since all \(\hat{X}_i\) operators commute when acting on distinct
qubits:

\[\hat{U}_B(\beta) = e^{-i\beta \sum_i \hat{X}_i} = \prod_{i=1}^{N} e^{-i\beta \hat{X}_i} = \prod_{i=1}^{N} R_X(2\beta)_i\]

where \(R_X(\theta) = e^{-i\theta \hat{X}/2}\).

\textbf{Full circuit:} For depth \(p\), the circuit is:

\begin{enumerate}
\def\labelenumi{\arabic{enumi}.}
\tightlist
\item
  Apply \(\hat{H}^{\otimes N}\) to \(|0\rangle^{\otimes N}\)
\item
  For \(k = 1, \ldots, p\):

  \begin{itemize}
  \tightlist
  \item
    Apply \(R_{ZZ}(-\gamma_k w_{ij})\) for each edge \((i,j) \in E\)
  \item
    Apply \(R_X(2\beta_k)\) to each qubit \(i \in V\)
  \end{itemize}
\item
  Measure in the computational basis
\end{enumerate}

\subsubsection{\texorpdfstring{2.4 The \(p \to \infty\) Limit and
Adiabatic
Connection}{2.4 The p \textbackslash to \textbackslash infty Limit and Adiabatic Connection}}\label{the-p-to-infty-limit-and-adiabatic-connection}

QAOA is formally equivalent to a digitised (Trotterised) quantum
annealing schedule. Consider the time-dependent Hamiltonian:

\[\hat{H}(t) = \left(1 - \frac{t}{T}\right)\hat{H}_B + \frac{t}{T}\hat{H}_C, \qquad t \in [0, T]\]

The adiabatic theorem guarantees that if \(T\) is large enough
(specifically \(T \sim O(1/\Delta_{\min}^2)\) where \(\Delta_{\min}\) is
the minimum spectral gap), the system tracks the instantaneous ground
state and ends in the ground state of \(\hat{H}_C\).

Trotterise this evolution into \(p\) steps with \(\Delta t = T/p\):

\[\hat{U} \approx \prod_{k=1}^{p} e^{-i\Delta t (1 - t_k/T)\hat{H}_B} \, e^{-i\Delta t (t_k/T)\hat{H}_C}\]

This is exactly the QAOA unitary with:

\[\beta_k = \Delta t\left(1 - \frac{k}{p}\right), \qquad \gamma_k = \Delta t \cdot \frac{k}{p}\]

So for fixed \(T\), as \(p \to \infty\), QAOA reproduces the adiabatic
evolution exactly (up to Trotter error \(O(\Delta t^2)\)). Since QAOA
optimises over all \((\boldsymbol{\gamma}, \boldsymbol{\beta})\) and is
not restricted to linear annealing schedules, it is strictly more
powerful than Trotterised adiabatic evolution: it can exploit shortcuts
to adiabaticity and counterdiabatic terms implicitly.

\subsubsection{\texorpdfstring{2.5 Analytical Expectation Value for
\(p=1\) on
\(C_4\)}{2.5 Analytical Expectation Value for p=1 on C\_4}}\label{analytical-expectation-value-for-p1-on-c_4}

For the cycle graph \(C_4\) with
\(\hat{H}_C = \sum_{(i,j) \in E} \frac{1}{2}(\hat{I} - \hat{Z}_i \hat{Z}_j)\)
(4 edges, 4 qubits), the \(p=1\) QAOA expectation value can be computed
analytically.

Start with \(|+\rangle^{\otimes 4}\). Apply
\(\hat{U}_C(\gamma) = \prod_{(i,j)} e^{i\gamma/2 \hat{Z}_i\hat{Z}_j}\)
(dropping global phases), then
\(\hat{U}_B(\beta) = \prod_i R_X(2\beta)_i\).

By the linearity of \(\langle \hat{H}_C \rangle\), we need:

\[\langle \hat{H}_C \rangle = \sum_{(i,j) \in E} \frac{1}{2}\bigl(1 - \langle \hat{Z}_i \hat{Z}_j \rangle\bigr)\]

By the symmetry of the cycle graph, all edge expectation values are
identical. Define
\(f(\gamma, \beta) = \langle \hat{Z}_0 \hat{Z}_1 \rangle\). Then
\(\langle \hat{H}_C \rangle = 4 \cdot \frac{1}{2}(1 - f) = 2(1 - f)\).

For the cycle on 4 qubits (ring topology), the \(p=1\) expectation value
is:

\[f(\gamma, \beta) = \langle +|^{\otimes 4}\, \hat{U}_C^\dagger(\gamma)\, \hat{U}_B^\dagger(\beta)\, \hat{Z}_0 \hat{Z}_1\, \hat{U}_B(\beta)\, \hat{U}_C(\gamma)\, |+\rangle^{\otimes 4}\]

Using the fact that
\(R_X(2\beta)^\dagger \hat{Z} R_X(2\beta) = \cos(2\beta)\hat{Z} - \sin(2\beta)\hat{Y}\)
and propagating through the \(ZZ\) rotations, one obtains (after
considerable algebra --- see Farhi et al.~Appendix or Wang et al., PRX
2018):

\[f(\gamma,\beta) = \cos^2(2\beta)\bigl[\cos(2\gamma) - \sin^2(2\gamma)\cos(2\gamma)\bigr] + \frac{1}{2}\sin^2(2\beta)\sin^2(2\gamma)\bigl[1 + \cos^2(2\gamma)\bigr]\sin(4\beta)(\ldots)\]

The full closed-form for a \(d\)-regular graph at \(p=1\) is given by
the Farhi-Goldstone-Gutmann formula. For the \(d\)-regular case with
edge \((i,j)\) where vertex \(i\) has \(d_i - 1\) other neighbours and
vertex \(j\) has \(d_j - 1\) other neighbours:

\[\langle \hat{Z}_i \hat{Z}_j \rangle = \frac{1}{2}\sin(4\beta)\sin(2\gamma)\bigl[\cos^{d_i-1}(2\gamma) + \cos^{d_j-1}(2\gamma)\bigr] - \frac{1}{4}\sin^2(2\beta)\sin^2(2\gamma)\Bigl[\cos^{n_{ij}}(2\gamma)\Bigr]\]

where \(n_{ij} = d_i + d_j - 2 - |N(i) \cap N(j)|\) accounts for shared
neighbours. For the 2-regular cycle \(C_4\), \(d_i = d_j = 2\), and
\(|N(0) \cap N(1)| = 0\), so \(n_{01} = 2\). Substituting:

\[\langle \hat{Z}_0\hat{Z}_1\rangle_{C_4} = \frac{1}{2}\sin(4\beta)\sin(2\gamma)\bigl[2\cos(2\gamma)\bigr] - \frac{1}{4}\sin^2(2\beta)\sin^2(2\gamma)\cos^2(2\gamma)\]

\[= \sin(4\beta)\sin(2\gamma)\cos(2\gamma) - \frac{1}{4}\sin^2(2\beta)\sin^2(2\gamma)\cos^2(2\gamma)\]

Then:

\[\langle \hat{H}_C \rangle = 2\bigl(1 - f(\gamma,\beta)\bigr) = 2 - 2\Bigl[\sin(4\beta)\sin(4\gamma)/2 - \frac{1}{4}\sin^2(2\beta)\sin^2(4\gamma)/4\Bigr]\]

The maximum of \(\langle \hat{H}_C \rangle\) over \((\gamma, \beta)\)
gives the QAOA \(p=1\) approximation ratio
\(r_1 = \langle \hat{H}_C \rangle_{\max} / C_{\max}\), where
\(C_{\max} = 4\).

Numerically, for \(C_4\): \(r_1 \approx 0.75\) (3 out of 4 edges),
achieved near \(\gamma^* \approx \pi/8\), \(\beta^* \approx \pi/8\).

\begin{center}\rule{0.5\linewidth}{0.5pt}\end{center}

\subsection{\texorpdfstring{3. The \((\beta, \gamma)\) Parameter
Landscape}{3. The (\textbackslash beta, \textbackslash gamma) Parameter Landscape}}\label{the-beta-gamma-parameter-landscape}

\subsubsection{3.1 Periodicity}\label{periodicity}

The QAOA landscape is periodic in both parameters. The periods depend on
the Hamiltonian structure:

\textbf{\(\gamma\)-periodicity:} The cost unitary
\(e^{-i\gamma \hat{H}_C}\) is periodic in \(\gamma\) with period
determined by the spectral structure of \(\hat{H}_C\). For MaxCut with
integer weights, the eigenvalues of \(\hat{H}_C\) are integers (or
half-integers), and the cost unitary is \(2\pi\)-periodic. More
precisely, if every eigenvalue difference of \(\hat{H}_C\) is an
integer, the period in \(\gamma\) is \(2\pi\). For the standard MaxCut
formulation \(\hat{H}_C = \sum_{(i,j)} \frac{1}{2}(I - Z_iZ_j)\), the
eigenvalue differences are integers, so:

\[\text{Period in } \gamma = 2\pi\]

However, exploiting the parity symmetry
\(\hat{H}_C(-\gamma) = \hat{H}_C(\gamma)\) and the fact that the
landscape has a reflection symmetry, it suffices to search
\(\gamma \in [0, \pi]\).

\textbf{\(\beta\)-periodicity:} The mixer unitary
\(e^{-i\beta \sum_i X_i}\): since \(X_i\) has eigenvalues \(\pm 1\), the
\(R_X(2\beta)\) gate is \(2\pi\)-periodic in \(2\beta\), i.e.,
\(\pi\)-periodic in \(\beta\):

\[\text{Period in } \beta = \pi\]

Again, the symmetry
\(\langle H_C \rangle(\beta) = \langle H_C \rangle(-\beta)\) allows
restriction to \(\beta \in [0, \pi/2]\).

\textbf{Combined search space for \(p=1\):}
\((\gamma, \beta) \in [0, \pi] \times [0, \pi/2]\) (or further reducible
depending on graph symmetries).

\subsubsection{3.2 Landscape Structure}\label{landscape-structure}

\textbf{Concentration of parameters (Brandao et al., 2018; Wurtz \&
Love, 2021):}

For \(d\)-regular graphs in the large \(N\) limit, the optimal QAOA
parameters \((\boldsymbol{\gamma}^*, \boldsymbol{\beta}^*)\) at fixed
\(p\) concentrate --- they become independent of the specific graph
instance and depend only on \(d\) and \(p\). This means:

\begin{itemize}
\tightlist
\item
  One can precompute ``universal'' optimal parameters on small instances
  and transfer them to larger graphs of the same regularity class.
\item
  The variance of the optimal energy across random \(d\)-regular graphs
  decays as \(O(1/N)\).
\end{itemize}

\textbf{Symmetry properties:}

\begin{enumerate}
\def\labelenumi{\arabic{enumi}.}
\tightlist
\item
  \textbf{Inversion symmetry:}
  \(\langle H_C \rangle(\boldsymbol{\gamma}, \boldsymbol{\beta}) = \langle H_C \rangle(-\boldsymbol{\gamma}, -\boldsymbol{\beta})\)
  --- the landscape is symmetric under simultaneous negation.
\item
  \textbf{Time-reversal:} For real Hamiltonians,
  \(\langle H_C \rangle(\boldsymbol{\gamma}, \boldsymbol{\beta}) = \langle H_C \rangle(\boldsymbol{\gamma}^R, \boldsymbol{\beta}^R)\)
  where \(R\) denotes reversal of parameter ordering.
\item
  \textbf{Graph automorphisms} induce symmetries in the landscape.
\end{enumerate}

\textbf{Good initial parameter strategies:}

\begin{itemize}
\tightlist
\item
  \textbf{INTERP (Zhou et al., 2020):} Optimise at depth \(p\), then
  linearly interpolate the parameters to initialise depth \(p+1\).
  Specifically, for \(p \to p+1\):
\end{itemize}

\[\gamma_k^{(p+1)} = \frac{k-1}{p}\gamma_{k-1}^{(p)} + \frac{p+1-k}{p}\gamma_k^{(p)}, \quad k = 1, \ldots, p+1\]

(similarly for \(\beta\)). This exploits the smoothness of optimal
parameter schedules.

\begin{itemize}
\tightlist
\item
  \textbf{FOURIER (Zhou et al., 2020):} Parameterise the QAOA angles in
  a Fourier basis:
\end{itemize}

\[\gamma_k = \sum_{q=1}^{Q} u_q \sin\!\left(\frac{(2k-1)q\pi}{2p}\right), \qquad \beta_k = \sum_{q=1}^{Q} v_q \cos\!\left(\frac{(2k-1)q\pi}{2p}\right)\]

where \(Q \leq p\) is the number of Fourier components. Optimisation
over \((u_1, \ldots, u_Q, v_1, \ldots, v_Q)\) is in a reduced
\(2Q\)-dimensional space instead of \(2p\), avoiding local minima.

\begin{itemize}
\tightlist
\item
  \textbf{TQA (Trotterised Quantum Annealing):} Use the linear schedule
  from the adiabatic connection (Section 2.4) as initialisation.
\end{itemize}

\subsubsection{3.3 Barren Plateaus}\label{barren-plateaus}

\textbf{Formal definition:} A variational circuit exhibits a barren
plateau if the variance of the cost function partial derivatives
vanishes exponentially with system size:

\[\text{Var}\!\left[\frac{\partial \langle \hat{H}_C \rangle}{\partial \theta_k}\right] \leq F(N)\]

where \(F(N) \in O(1/b^N)\) for some \(b > 1\).

For QAOA specifically, the situation is nuanced:

\begin{enumerate}
\def\labelenumi{\arabic{enumi}.}
\item
  \textbf{Hardware-efficient random ansätze} (McClean et al., 2018): For
  circuits that form approximate 2-designs, the gradient variance decays
  as \(O(2^{-N})\), guaranteeing barren plateaus.
\item
  \textbf{QAOA with fixed \(p\):} QAOA at constant depth \(p\) does NOT
  suffer from barren plateaus in the same way because the circuit is
  highly structured (alternating cost and mixer unitaries) and shallow.
  The gradient variance scales polynomially (often as
  \(\Theta(1/\text{poly}(N))\)), which is trainable.
\item
  \textbf{QAOA with \(p = O(\text{poly}(N))\):} As \(p\) grows
  polynomially with \(N\), the landscape can develop barren plateau
  features depending on the cost function locality. For local cost
  functions (each term acts on \(O(1)\) qubits), Cerezo et al.~(2021)
  showed that shallow circuits have well-behaved gradients, while deep
  circuits suffer barren plateaus regardless.
\end{enumerate}

\textbf{Practical consequence:} For the typical QAOA regime
(\(p \leq 10\), \(N \leq 30\)), barren plateaus are not the primary
obstacle. Instead, the proliferation of local minima and the
non-convexity of the landscape are the main challenges.

\subsubsection{3.4 Phase Transition Behaviour with Local Field
Strength}\label{phase-transition-behaviour-with-local-field-strength}

The 2025 result in \emph{Quantum} (Boulebnane \& Montanaro, 2025; also
Kremenetski et al., 2025) studies the QAOA landscape on
anti-ferromagnetic Ising models with longitudinal fields:

\[\hat{H}_C = -\sum_{(i,j)} J_{ij}\hat{Z}_i\hat{Z}_j + h\sum_i \hat{Z}_i\]

where \(J_{ij} < 0\) (anti-ferromagnetic). Three distinct regimes are
identified as \(h\) varies:

\begin{enumerate}
\def\labelenumi{\arabic{enumi}.}
\item
  \textbf{Trivial regime} (\(|h| \gg |J|\)): The ground state is the
  fully polarised state (all spins aligned with the field). QAOA finds
  this trivially even at \(p=1\) --- the landscape has a single dominant
  minimum.
\item
  \textbf{Anti-ferromagnetic regime} (\(|h| \ll |J|\)): The ground state
  exhibits Néel-type anti-ferromagnetic ordering. The landscape at low
  \(p\) is highly structured with multiple degenerate minima
  corresponding to the broken symmetry. The QAOA performance depends
  strongly on the graph geometry (bipartite vs.~frustrated).
\item
  \textbf{Critical regime} (\(|h| \sim |J|\), near the quantum phase
  transition): The spectral gap \(\Delta\) closes polynomially or
  logarithmically with \(N\), the landscape becomes complex with many
  nearly degenerate minima, and the approximation ratio achievable at
  fixed \(p\) drops sharply. This is where QAOA most struggles --- and
  where the landscape structure is richest from a visualisation
  perspective.
\end{enumerate}

The key finding is that QAOA performance undergoes a sharp transition at
the critical field strength, and the landscape topology (number of local
minima, gradient magnitudes, Hessian structure) changes qualitatively
across the phase boundary.

\begin{center}\rule{0.5\linewidth}{0.5pt}\end{center}

\subsection{4. Classical Optimization of QAOA
Parameters}\label{classical-optimization-of-qaoa-parameters}

\subsubsection{4.1 Parameter-Shift Rule: Full
Derivation}\label{parameter-shift-rule-full-derivation}

For a parameterised unitary of the form \(U(\theta) = e^{-i\theta G/2}\)
where \(G\) is a generator with eigenvalues \(\pm r\) (typically
\(r = 1\) for Pauli generators), the parameter-shift rule states:

\[\frac{\partial}{\partial \theta}\langle \psi(\theta) | \hat{O} | \psi(\theta)\rangle = \frac{r}{2}\Bigl[\langle \hat{O}\rangle_{\theta + \pi/(2r)} - \langle \hat{O}\rangle_{\theta - \pi/(2r)}\Bigr]\]

\textbf{Derivation:} Let
\(|\psi(\theta)\rangle = U(\theta)|\psi_0\rangle\) with
\(U(\theta) = e^{-i\theta G/2}\). Decompose \(|\psi_0\rangle\) in the
eigenbasis of \(G\):
\(G = r|+\rangle\langle +| - r|-\rangle\langle -|\).

\[|\psi_0\rangle = c_+|+\rangle + c_-|-\rangle\]

\[|\psi(\theta)\rangle = c_+ e^{-ir\theta/2}|+\rangle + c_- e^{ir\theta/2}|-\rangle\]

Let \(f(\theta) = \langle\psi(\theta)|\hat{O}|\psi(\theta)\rangle\).
Computing:

\[f(\theta) = |c_+|^2 O_{++} + |c_-|^2 O_{--} + c_+^* c_- e^{ir\theta} O_{+-} + c_- ^* c_+ e^{-ir\theta} O_{-+}\]

where \(O_{ab} = \langle a|\hat{O}|b\rangle\). This is a sinusoidal
function with frequency \(r\):

\[f(\theta) = A + B\cos(r\theta) + C\sin(r\theta)\]

Therefore:

\[f'(\theta) = -Br\sin(r\theta) + Cr\cos(r\theta) = r\bigl[f(\theta + \pi/(2r)) - f(\theta - \pi/(2r))\bigr] / 2\]

which follows from the identity:

\[\frac{f(\theta + s) - f(\theta - s)}{2} = B\bigl[\cos r(\theta+s) - \cos r(\theta-s)\bigr]/2 + C\bigl[\sin r(\theta + s) - \sin r(\theta-s)\bigr]/2\]

Setting \(s = \pi/(2r)\):

\[= -B\sin(r\theta)\sin(\pi/2) + C\cos(r\theta)\sin(\pi/2) = -B\sin(r\theta) + C\cos(r\theta) = f'(\theta)/r\]

Hence \(f'(\theta) = r[f(\theta + \pi/(2r)) - f(\theta - \pi/(2r))]/2\).
For standard Pauli generators, \(r = 1\) and the shift is \(\pm \pi/2\):

\[\frac{\partial \langle \hat{O} \rangle}{\partial \theta} = \frac{\langle \hat{O}\rangle_{\theta + \pi/2} - \langle \hat{O}\rangle_{\theta - \pi/2}}{2}\]

\textbf{For QAOA:} The cost unitary parameters \(\gamma_k\) multiply
\(\hat{H}_C\), which has a richer spectrum than \(\pm 1\). The
\(ZZ\)-rotation gates \(e^{-i\gamma w_{ij} \hat{Z}_i \hat{Z}_j/2}\)
individually satisfy the two-eigenvalue condition with \(r = w_{ij}\),
so the parameter-shift rule applies gate-by-gate. However, the full
circuit expectation value as a function of \(\gamma_k\) involves
multiple \(ZZ\)-gates at the same \(\gamma_k\), leading to a
multi-frequency sinusoidal. In practice, one can either: - Decompose the
cost layer and compute gradients one gate at a time (more circuits). -
Use the general parameter-shift rule for multi-frequency functions
(Wierichs et al., 2022).

\subsubsection{4.2 Optimizer Comparison}\label{optimizer-comparison}

{\def\LTcaptype{none} % do not increment counter
\begin{longtable}[]{@{}
  >{\raggedright\arraybackslash}p{(\linewidth - 8\tabcolsep) * \real{0.1642}}
  >{\raggedright\arraybackslash}p{(\linewidth - 8\tabcolsep) * \real{0.0896}}
  >{\raggedright\arraybackslash}p{(\linewidth - 8\tabcolsep) * \real{0.2239}}
  >{\raggedright\arraybackslash}p{(\linewidth - 8\tabcolsep) * \real{0.2537}}
  >{\raggedright\arraybackslash}p{(\linewidth - 8\tabcolsep) * \real{0.2687}}@{}}
\toprule\noalign{}
\begin{minipage}[b]{\linewidth}\raggedright
Optimizer
\end{minipage} & \begin{minipage}[b]{\linewidth}\raggedright
Type
\end{minipage} & \begin{minipage}[b]{\linewidth}\raggedright
Gradient-free?
\end{minipage} & \begin{minipage}[b]{\linewidth}\raggedright
Noisy-tolerant?
\end{minipage} & \begin{minipage}[b]{\linewidth}\raggedright
Typical use case
\end{minipage} \\
\midrule\noalign{}
\endhead
\bottomrule\noalign{}
\endlastfoot
\textbf{COBYLA} & Trust-region, gradient-free & Yes & Moderate & Small
\(p\), few parameters, hardware runs \\
\textbf{L-BFGS-B} & Quasi-Newton & No (needs gradient) & Poor &
Noiseless simulation, smooth landscapes \\
\textbf{SPSA} & Stochastic gradient approximation & Effectively yes (2
evals/step) & \textbf{Excellent} & Hardware, noisy estimates, many
parameters \\
\textbf{Adam} & Momentum-based gradient descent & No (needs gradient) &
Moderate & Differentiable simulators (autograd) \\
\end{longtable}
}

\textbf{COBYLA} (Constrained Optimization BY Linear Approximations):
Uses a simplex-like trust region without gradients. Requires \(O(d)\)
function evaluations per step where \(d\) is the parameter dimension.
Works well for \(d \leq 20\) but scales poorly. On noisy hardware, the
trust region can be destabilised by shot noise.

\textbf{L-BFGS-B} (Limited-memory BFGS with Box constraints):
Approximates the Hessian from gradient history. Converges superlinearly
near optima in noiseless settings. Completely unusable on hardware
because noise in the gradient estimates corrupts the Hessian
approximation, causing divergence.

\textbf{SPSA} (Simultaneous Perturbation Stochastic Approximation):
Estimates the gradient using only 2 function evaluations (regardless of
dimension \(d\)) by perturbing all parameters simultaneously with random
\(\pm c_k\) shifts:

\[\hat{g}_k = \frac{f(\boldsymbol{\theta} + c_k \boldsymbol{\Delta}) - f(\boldsymbol{\theta} - c_k \boldsymbol{\Delta})}{2c_k} \boldsymbol{\Delta}^{-1}\]

where \(\boldsymbol{\Delta}\) is a random perturbation vector
(Rademacher). The gain schedule \(a_k, c_k\) must be tuned. SPSA is the
go-to optimiser for hardware because it needs constant circuit
evaluations per step and is inherently noise-tolerant.

\textbf{Adam:} Adaptive moment estimation with bias correction. Works
brilliantly with exact gradients from autograd simulators (PennyLane's
backpropagation or adjoint differentiation). Not directly usable on
hardware unless combined with parameter-shift gradients, which adds
overhead.

\subsubsection{4.3 Gradient-Free vs Gradient-Based on Noisy
Hardware}\label{gradient-free-vs-gradient-based-on-noisy-hardware}

The fundamental tension:

\begin{itemize}
\item
  \textbf{Gradient-based} methods (parameter-shift + Adam/L-BFGS-B) give
  accurate search directions but require \(O(d)\) circuit evaluations
  per gradient (2 per parameter via parameter-shift), each of which has
  shot noise. The total gradient estimation error scales as
  \(O(d/\sqrt{S})\) where \(S\) is the number of shots.
\item
  \textbf{Gradient-free} methods (COBYLA, Nelder-Mead, SPSA) avoid the
  \(O(d)\) overhead but rely on noisy function evaluations to infer
  descent directions, which may require more iterations to converge.
\end{itemize}

\textbf{Practical guidance:} - For simulators: use autograd (PennyLane's
\texttt{diff\_method="backprop"}) with Adam. This is exact and fast. -
For hardware with \(d \leq 10\) (QAOA \(p \leq 5\)): COBYLA or SPSA. -
For hardware with \(d > 10\): SPSA is the only viable option due to its
constant evaluation cost per step. - Hybrid: optimise on simulator to
find a good initialisation, then fine-tune on hardware with SPSA.

\begin{center}\rule{0.5\linewidth}{0.5pt}\end{center}

\section{Part 2 --- Software
Architecture}\label{part-2-software-architecture}

\subsection{Project Structure}\label{project-structure}

\begin{verbatim}
qaoa_explorer/
├── main.py                  # Orchestration, CLI, entry point
├── graph_generator.py       # Graph construction, adjacency → Ising
├── qaoa_circuit.py          # QAOA circuit construction (PennyLane)
├── optimizer.py             # Classical optimization loop
├── landscape_sampler.py     # Grid sampling of (β, γ) space
├── noise_model.py           # Noise simulation (PennyLane + QuTiP)
├── hardware_runner.py       # IBM Quantum / IQM hardware submission
├── error_mitigation.py      # ZNE implementation
├── renderer.py              # OpenGL 3D landscape visualization
├── requirements.txt
├── config.yaml              # Default configuration
└── tests/
    ├── test_graph.py
    ├── test_circuit.py
    ├── test_optimizer.py
    ├── test_sampler.py
    └── test_renderer.py
\end{verbatim}

\subsection{Data Flow Diagram}\label{data-flow-diagram}

\begin{verbatim}
                     ┌───────────────┐
                     │  config.yaml  │
                     └──────┬────────┘
                            │
                     ┌──────▼────────┐
                     │   main.py     │  CLI / orchestration
                     └──────┬────────┘
                            │
              ┌─────────────┼─────────────────┐
              │             │                  │
     ┌────────▼──────┐ ┌───▼────────────┐ ┌───▼──────────────┐
     │graph_generator│ │landscape_sampler│ │   optimizer.py   │
     │   .py         │ │   .py           │ │                  │
     └────────┬──────┘ └───┬────────────┘ └───┬──────────────┘
              │            │                   │
              │   ┌────────▼─────────┐         │
              └──►│ qaoa_circuit.py  │◄────────┘
                  └────────┬─────────┘
                           │
           ┌───────────────┼───────────────────┐
           │               │                    │
  ┌────────▼──────┐ ┌──────▼───────┐  ┌────────▼──────────┐
  │noise_model.py │ │hardware_     │  │error_mitigation.py│
  │(default.mixed)│ │runner.py     │  │  (ZNE)            │
  └────────┬──────┘ │(IBM/IQM)     │  └────────┬──────────┘
           │        └──────┬───────┘           │
           └───────────────┼───────────────────┘
                           │
                    ┌──────▼───────┐
                    │ renderer.py  │  OpenGL surface
                    │ (β,γ,⟨H_C⟩) │
                    └──────────────┘
\end{verbatim}

\subsection{Module Specifications}\label{module-specifications}

\subsubsection{\texorpdfstring{1.
\texttt{graph\_generator.py}}{1. graph\_generator.py}}\label{graph_generator.py}

\textbf{Purpose:} Construct graphs and convert adjacency representations
to Ising Hamiltonians.

\begin{Shaded}
\begin{Highlighting}[]
\ImportTok{from}\NormalTok{ \_\_future\_\_ }\ImportTok{import}\NormalTok{ annotations}
\ImportTok{import}\NormalTok{ networkx }\ImportTok{as}\NormalTok{ nx}
\ImportTok{import}\NormalTok{ pennylane }\ImportTok{as}\NormalTok{ qml}


\KeywordTok{def}\NormalTok{ create\_graph(}
\NormalTok{    graph\_type: }\BuiltInTok{str}\NormalTok{,}
\NormalTok{    n\_nodes: }\BuiltInTok{int}\NormalTok{,}
\NormalTok{    seed: }\BuiltInTok{int} \OperatorTok{|} \VariableTok{None} \OperatorTok{=} \VariableTok{None}\NormalTok{,}
    \OperatorTok{**}\NormalTok{kwargs,}
\NormalTok{) }\OperatorTok{{-}\textgreater{}}\NormalTok{ nx.Graph:}
    \CommentTok{"""Create a NetworkX graph of the specified type.}

\CommentTok{    Supported types: \textquotesingle{}cycle\textquotesingle{}, \textquotesingle{}complete\textquotesingle{}, \textquotesingle{}erdos\_renyi\textquotesingle{}, \textquotesingle{}random\_regular\textquotesingle{},}
\CommentTok{    \textquotesingle{}grid\_2d\textquotesingle{}, \textquotesingle{}petersen\textquotesingle{}, \textquotesingle{}custom\textquotesingle{}. Stochastic families accept a random}
\CommentTok{    seed for reproducibility. Extra kwargs control family{-}specific}
\CommentTok{    parameters (e.g. p=0.5 for Erdos{-}Renyi, d=3 for random regular).}

\CommentTok{    Returns:}
\CommentTok{        A weighted undirected NetworkX graph with unit{-}weight edges}
\CommentTok{        by default.}
\CommentTok{    """}
\NormalTok{    ...}


\KeywordTok{def}\NormalTok{ adjacency\_to\_ising(}
\NormalTok{    graph: nx.Graph,}
\NormalTok{) }\OperatorTok{{-}\textgreater{}} \BuiltInTok{tuple}\NormalTok{[}\BuiltInTok{list}\NormalTok{[}\BuiltInTok{float}\NormalTok{], }\BuiltInTok{list}\NormalTok{[}\BuiltInTok{tuple}\NormalTok{[}\BuiltInTok{int}\NormalTok{, }\BuiltInTok{int}\NormalTok{]], }\BuiltInTok{float}\NormalTok{]:}
    \CommentTok{"""Convert a graph\textquotesingle{}s adjacency structure to Ising Hamiltonian coefficients.}

\CommentTok{    Uses the MaxCut formulation:}
\CommentTok{        H\_C = sum\_\{(i,j)\} w\_ij/2 (I {-} Z\_i Z\_j)}

\CommentTok{    For each edge (i,j) with weight w\_ij, produces coefficient {-}w\_ij/2}
\CommentTok{    paired with the Pauli{-}Z index pair (i,j), plus a constant energy}
\CommentTok{    offset sum\_\{(i,j)\} w\_ij/2.}

\CommentTok{    Returns:}
\CommentTok{        Tuple of (coefficients, observable\_index\_pairs, offset).}
\CommentTok{    """}
\NormalTok{    ...}


\KeywordTok{def}\NormalTok{ build\_cost\_hamiltonian(graph: nx.Graph) }\OperatorTok{{-}\textgreater{}}\NormalTok{ qml.Hamiltonian:}
    \CommentTok{"""Wrap the Ising conversion into a PennyLane Hamiltonian.}

\CommentTok{    Each edge contributes a {-}w\_ij/2 * Z\_i @ Z\_j term. The identity}
\CommentTok{    offset is included as an explicit qml.Identity term so that}
\CommentTok{    absolute energy values are preserved.}

\CommentTok{    Returns:}
\CommentTok{        PennyLane Hamiltonian ready for circuit evaluation.}
\CommentTok{    """}
\NormalTok{    ...}
\end{Highlighting}
\end{Shaded}

\emph{(Full implementation:
\texttt{report1\_scripts/graph\_generator.py})}

\subsubsection{\texorpdfstring{2.
\texttt{qaoa\_circuit.py}}{2. qaoa\_circuit.py}}\label{qaoa_circuit.py}

\textbf{Purpose:} Construct and evaluate QAOA circuits using PennyLane.

\begin{Shaded}
\begin{Highlighting}[]
\ImportTok{from}\NormalTok{ \_\_future\_\_ }\ImportTok{import}\NormalTok{ annotations}
\ImportTok{import}\NormalTok{ numpy }\ImportTok{as}\NormalTok{ np}
\ImportTok{import}\NormalTok{ pennylane }\ImportTok{as}\NormalTok{ qml}


\KeywordTok{def}\NormalTok{ cost\_layer(}
\NormalTok{    gamma: }\BuiltInTok{float}\NormalTok{,}
\NormalTok{    graph\_edges: }\BuiltInTok{list}\NormalTok{[}\BuiltInTok{tuple}\NormalTok{[}\BuiltInTok{int}\NormalTok{, }\BuiltInTok{int}\NormalTok{, }\BuiltInTok{float}\NormalTok{]],}
\NormalTok{) }\OperatorTok{{-}\textgreater{}} \VariableTok{None}\NormalTok{:}
    \CommentTok{"""Apply the cost unitary U\_C(gamma) = exp({-}i gamma H\_C).}

\CommentTok{    Decomposes each weighted edge (i, j, w\_ij) into an IsingZZ rotation:}
\CommentTok{        exp({-}i gamma w\_ij Z\_i Z\_j / 2)}
\CommentTok{    implemented via qml.IsingZZ(gamma * w\_ij, wires=[i, j]).}
\CommentTok{    All ZZ terms commute, so ordering is irrelevant.}
\CommentTok{    """}
\NormalTok{    ...}


\KeywordTok{def}\NormalTok{ mixer\_layer(beta: }\BuiltInTok{float}\NormalTok{, n\_qubits: }\BuiltInTok{int}\NormalTok{) }\OperatorTok{{-}\textgreater{}} \VariableTok{None}\NormalTok{:}
    \CommentTok{"""Apply the mixer unitary U\_B(beta) = exp({-}i beta H\_B).}

\CommentTok{    Applies independent single{-}qubit RX(2*beta) rotations on every}
\CommentTok{    qubit, since H\_B = sum\_i X\_i factorises across qubits.}
\CommentTok{    """}
\NormalTok{    ...}


\KeywordTok{def}\NormalTok{ qaoa\_ansatz(}
\NormalTok{    params: np.ndarray,}
\NormalTok{    cost\_hamiltonian: qml.Hamiltonian,}
\NormalTok{    n\_qubits: }\BuiltInTok{int}\NormalTok{,}
\NormalTok{    p: }\BuiltInTok{int}\NormalTok{,}
\NormalTok{) }\OperatorTok{{-}\textgreater{}} \VariableTok{None}\NormalTok{:}
    \CommentTok{"""Apply the full QAOA state{-}preparation sequence inside a QNode context.}

\CommentTok{    Prepares |+\textgreater{}\^{}N via a Hadamard layer, then applies p alternating}
\CommentTok{    rounds of cost and mixer unitaries:}
\CommentTok{        prod\_\{k=1\}\^{}\{p\} [U\_B(beta\_k) U\_C(gamma\_k)]}

\CommentTok{    Args:}
\CommentTok{        params: Array of shape (2p,). First p entries are gamma\_k,}
\CommentTok{                last p entries are beta\_k.}
\CommentTok{    """}
\NormalTok{    ...}


\KeywordTok{def}\NormalTok{ build\_qaoa\_circuit(}
\NormalTok{    cost\_hamiltonian: qml.Hamiltonian,}
\NormalTok{    n\_qubits: }\BuiltInTok{int}\NormalTok{,}
\NormalTok{    p: }\BuiltInTok{int}\NormalTok{,}
\NormalTok{    device: qml.Device,}
\NormalTok{) }\OperatorTok{{-}\textgreater{}}\NormalTok{ qml.QNode:}
    \CommentTok{"""Wrap the QAOA ansatz into a PennyLane QNode bound to the given device.}

\CommentTok{    The QNode accepts a parameter array of shape (2p,) {-}{-} the first p}
\CommentTok{    entries are the gamma\_k values, the last p are the beta\_k values {-}{-}}
\CommentTok{    and returns \textless{}H\_C\textgreater{}. Differentiation method is set to \textquotesingle{}best\textquotesingle{}}
\CommentTok{    (backprop for simulators, parameter{-}shift for hardware).}

\CommentTok{    Returns:}
\CommentTok{        A PennyLane QNode returning the cost expectation value.}
\CommentTok{    """}
\NormalTok{    ...}
\end{Highlighting}
\end{Shaded}

\emph{(Full implementation: \texttt{report1\_scripts/qaoa\_circuit.py})}

\subsubsection{\texorpdfstring{3.
\texttt{optimizer.py}}{3. optimizer.py}}\label{optimizer.py}

\textbf{Purpose:} Classical optimization of QAOA parameters with
multiple backends.

\begin{Shaded}
\begin{Highlighting}[]
\ImportTok{from}\NormalTok{ \_\_future\_\_ }\ImportTok{import}\NormalTok{ annotations}
\ImportTok{from}\NormalTok{ dataclasses }\ImportTok{import}\NormalTok{ dataclass, field}
\ImportTok{from}\NormalTok{ typing }\ImportTok{import}\NormalTok{ Callable}
\ImportTok{import}\NormalTok{ numpy }\ImportTok{as}\NormalTok{ np}
\ImportTok{import}\NormalTok{ pennylane }\ImportTok{as}\NormalTok{ qml}


\AttributeTok{@dataclass}
\KeywordTok{class}\NormalTok{ OptimizationResult:}
    \CommentTok{"""Container for QAOA optimization results.}

\CommentTok{    Attributes:}
\CommentTok{        best\_params: Optimal parameter vector of shape (2p,).}
\CommentTok{        optimal\_energy: Best \textless{}H\_C\textgreater{} value found.}
\CommentTok{        trajectory: List of (params, energy) pairs at each iteration.}
\CommentTok{        n\_evaluations: Total number of circuit evaluations.}
\CommentTok{        converged: Whether the optimizer reached the tolerance.}
\CommentTok{    """}
\NormalTok{    best\_params: np.ndarray}
\NormalTok{    optimal\_energy: }\BuiltInTok{float}
\NormalTok{    trajectory: }\BuiltInTok{list}\NormalTok{[}\BuiltInTok{tuple}\NormalTok{[np.ndarray, }\BuiltInTok{float}\NormalTok{]] }\OperatorTok{=}\NormalTok{ field(default\_factory}\OperatorTok{=}\BuiltInTok{list}\NormalTok{)}
\NormalTok{    n\_evaluations: }\BuiltInTok{int} \OperatorTok{=} \DecValTok{0}
\NormalTok{    converged: }\BuiltInTok{bool} \OperatorTok{=} \VariableTok{False}


\KeywordTok{def}\NormalTok{ optimize\_qaoa(}
\NormalTok{    qnode: qml.QNode,}
\NormalTok{    p: }\BuiltInTok{int}\NormalTok{,}
\NormalTok{    method: }\BuiltInTok{str} \OperatorTok{=} \StringTok{"COBYLA"}\NormalTok{,}
\NormalTok{    init\_params: np.ndarray }\OperatorTok{|} \VariableTok{None} \OperatorTok{=} \VariableTok{None}\NormalTok{,}
\NormalTok{    maxiter: }\BuiltInTok{int} \OperatorTok{=} \DecValTok{500}\NormalTok{,}
\NormalTok{    tol: }\BuiltInTok{float} \OperatorTok{=} \FloatTok{1e{-}6}\NormalTok{,}
\NormalTok{    callback: Callable[[np.ndarray, }\BuiltInTok{float}\NormalTok{, }\BuiltInTok{int}\NormalTok{], }\VariableTok{None}\NormalTok{] }\OperatorTok{|} \VariableTok{None} \OperatorTok{=} \VariableTok{None}\NormalTok{,}
\NormalTok{) }\OperatorTok{{-}\textgreater{}}\NormalTok{ OptimizationResult:}
    \CommentTok{"""Run classical optimisation of the 2p QAOA parameters.}

\CommentTok{    Wraps the QNode in a negated objective (QAOA maximises \textless{}H\_C\textgreater{} but}
\CommentTok{    standard optimisers minimise).}

\CommentTok{    Supported methods:}
\CommentTok{        {-} \textquotesingle{}COBYLA\textquotesingle{}: Gradient{-}free trust{-}region via scipy.optimize.minimize.}
\CommentTok{          Best for few parameters (d \textless{}= 20).}
\CommentTok{        {-} \textquotesingle{}L{-}BFGS{-}B\textquotesingle{}: Quasi{-}Newton requiring exact gradients via the}
\CommentTok{          parameter{-}shift rule. Fast near optima in noiseless settings.}
\CommentTok{        {-} \textquotesingle{}SPSA\textquotesingle{}: Stochastic perturbation gradient approximation. Uses}
\CommentTok{          only 2 function evaluations per step regardless of dimension,}
\CommentTok{          with Rademacher perturbation vectors and standard gain schedule}
\CommentTok{          a\_k = a0/(k+A)\^{}alpha, c\_k = c0/k\^{}gamma.}
\CommentTok{        {-} \textquotesingle{}Adam\textquotesingle{}: PennyLane\textquotesingle{}s adaptive moment{-}based gradient descent,}
\CommentTok{          suited for autograd{-}differentiated simulators.}

\CommentTok{    If init\_params is None, gamma\_k \textasciitilde{} Uniform(0, pi) and}
\CommentTok{    beta\_k \textasciitilde{} Uniform(0, pi/2). The optional callback is invoked at}
\CommentTok{    each iteration with (params, energy, iteration).}
\CommentTok{    """}
\NormalTok{    ...}


\KeywordTok{def}\NormalTok{ parameter\_shift\_gradient(}
\NormalTok{    qnode: qml.QNode,}
\NormalTok{    params: np.ndarray,}
\NormalTok{    shift: }\BuiltInTok{float} \OperatorTok{=}\NormalTok{ np.pi }\OperatorTok{/} \DecValTok{2}\NormalTok{,}
\NormalTok{) }\OperatorTok{{-}\textgreater{}}\NormalTok{ np.ndarray:}
    \CommentTok{"""Compute the gradient of the QNode via the parameter{-}shift rule.}

\CommentTok{    For each parameter theta\_i, evaluates at theta\_i +/{-} pi/2:}
\CommentTok{        d\textless{}H\_C\textgreater{}/d(theta\_i) = [f(theta\_i + pi/2) {-} f(theta\_i {-} pi/2)] / 2}

\CommentTok{    Requires 2d circuit evaluations for d parameters. Hardware{-}compatible}
\CommentTok{    (does not rely on autograd). Exact for gates whose generators have}
\CommentTok{    eigenvalues +/{-}1 (all Pauli{-}based rotations in the QAOA circuit).}

\CommentTok{    Returns:}
\CommentTok{        Gradient array of shape (2p,).}
\CommentTok{    """}
\NormalTok{    ...}
\end{Highlighting}
\end{Shaded}

\emph{(Full implementation: \texttt{report1\_scripts/optimizer.py})}

\subsubsection{\texorpdfstring{4.
\texttt{landscape\_sampler.py}}{4. landscape\_sampler.py}}\label{landscape_sampler.py}

\textbf{Purpose:} Efficient grid sampling of the \((\beta, \gamma)\)
landscape.

\begin{Shaded}
\begin{Highlighting}[]
\ImportTok{from}\NormalTok{ \_\_future\_\_ }\ImportTok{import}\NormalTok{ annotations}
\ImportTok{from}\NormalTok{ dataclasses }\ImportTok{import}\NormalTok{ dataclass}
\ImportTok{import}\NormalTok{ numpy }\ImportTok{as}\NormalTok{ np}
\ImportTok{import}\NormalTok{ pennylane }\ImportTok{as}\NormalTok{ qml}


\AttributeTok{@dataclass}
\KeywordTok{class}\NormalTok{ LandscapeData:}
    \CommentTok{"""Container for energy landscape sampling results.}

\CommentTok{    Attributes:}
\CommentTok{        gamma\_vals: 1D array of gamma sample points, shape (N\_gamma,).}
\CommentTok{        beta\_vals: 1D array of beta sample points, shape (N\_beta,).}
\CommentTok{        energies: 2D energy array of shape (N\_gamma, N\_beta).}
\CommentTok{        approx\_ratios: 2D approximation ratios normalised by C\_max,}
\CommentTok{                        shape (N\_gamma, N\_beta).}
\CommentTok{    """}
\NormalTok{    gamma\_vals: np.ndarray}
\NormalTok{    beta\_vals: np.ndarray}
\NormalTok{    energies: np.ndarray}
\NormalTok{    approx\_ratios: np.ndarray}


\KeywordTok{def}\NormalTok{ sample\_landscape(}
\NormalTok{    qnode: qml.QNode,}
\NormalTok{    p: }\BuiltInTok{int}\NormalTok{,}
\NormalTok{    gamma\_range: }\BuiltInTok{tuple}\NormalTok{[}\BuiltInTok{float}\NormalTok{, }\BuiltInTok{float}\NormalTok{],}
\NormalTok{    beta\_range: }\BuiltInTok{tuple}\NormalTok{[}\BuiltInTok{float}\NormalTok{, }\BuiltInTok{float}\NormalTok{],}
\NormalTok{    n\_gamma: }\BuiltInTok{int}\NormalTok{,}
\NormalTok{    n\_beta: }\BuiltInTok{int}\NormalTok{,}
\NormalTok{    fixed\_params: np.ndarray }\OperatorTok{|} \VariableTok{None} \OperatorTok{=} \VariableTok{None}\NormalTok{,}
\NormalTok{    c\_max: }\BuiltInTok{float} \OperatorTok{=} \FloatTok{1.0}\NormalTok{,}
\NormalTok{) }\OperatorTok{{-}\textgreater{}}\NormalTok{ LandscapeData:}
    \CommentTok{"""Evaluate \textless{}H\_C\textgreater{} on a uniform (N\_gamma x N\_beta) grid.}

\CommentTok{    For p=1 the two parameters are swept directly. For p\textgreater{}1, all}
\CommentTok{    parameters except the first gamma and first beta are held fixed}
\CommentTok{    at fixed\_params, producing a 2D slice through the higher{-}dimensional}
\CommentTok{    landscape.}

\CommentTok{    Returns:}
\CommentTok{        LandscapeData with energies and approximation ratios.}
\CommentTok{    """}
\NormalTok{    ...}


\KeywordTok{def}\NormalTok{ sample\_landscape\_parallel(}
\NormalTok{    qnode: qml.QNode,}
\NormalTok{    p: }\BuiltInTok{int}\NormalTok{,}
\NormalTok{    gamma\_range: }\BuiltInTok{tuple}\NormalTok{[}\BuiltInTok{float}\NormalTok{, }\BuiltInTok{float}\NormalTok{],}
\NormalTok{    beta\_range: }\BuiltInTok{tuple}\NormalTok{[}\BuiltInTok{float}\NormalTok{, }\BuiltInTok{float}\NormalTok{],}
\NormalTok{    n\_gamma: }\BuiltInTok{int}\NormalTok{,}
\NormalTok{    n\_beta: }\BuiltInTok{int}\NormalTok{,}
\NormalTok{    c\_max: }\BuiltInTok{float} \OperatorTok{=} \FloatTok{1.0}\NormalTok{,}
\NormalTok{) }\OperatorTok{{-}\textgreater{}}\NormalTok{ LandscapeData:}
    \CommentTok{"""Batched grid evaluation using qml.execute over pre{-}constructed tapes.}

\CommentTok{    Builds parameter grid via np.meshgrid (indexing=\textquotesingle{}ij\textquotesingle{}), flattens all}
\CommentTok{    (N\_gamma * N\_beta) parameter sets, executes in a single call to}
\CommentTok{    qml.execute, then reshapes into the 2D energy array. Avoids}
\CommentTok{    Python{-}loop overhead; significantly faster on simulators.}

\CommentTok{    Returns:}
\CommentTok{        LandscapeData with energies and approximation ratios.}
\CommentTok{    """}
\NormalTok{    ...}
\end{Highlighting}
\end{Shaded}

\emph{(Full implementation:
\texttt{report1\_scripts/landscape\_sampler.py})}

\subsubsection{\texorpdfstring{5.
\texttt{noise\_model.py}}{5. noise\_model.py}}\label{noise_model.py}

\textbf{Purpose:} Noise simulation using PennyLane's
\texttt{default.mixed} and QuTiP for Lindblad master equation analysis.

\begin{Shaded}
\begin{Highlighting}[]
\ImportTok{from}\NormalTok{ \_\_future\_\_ }\ImportTok{import}\NormalTok{ annotations}
\ImportTok{import}\NormalTok{ numpy }\ImportTok{as}\NormalTok{ np}
\ImportTok{import}\NormalTok{ pennylane }\ImportTok{as}\NormalTok{ qml}
\ImportTok{import}\NormalTok{ qutip}


\KeywordTok{def}\NormalTok{ create\_noisy\_device(}
\NormalTok{    n\_qubits: }\BuiltInTok{int}\NormalTok{,}
\NormalTok{    depolarizing\_rate: }\BuiltInTok{float} \OperatorTok{=} \FloatTok{0.001}\NormalTok{,}
\NormalTok{    thermal\_relaxation\_t1: }\BuiltInTok{float} \OperatorTok{=} \FloatTok{50e{-}6}\NormalTok{,}
\NormalTok{    thermal\_relaxation\_t2: }\BuiltInTok{float} \OperatorTok{=} \FloatTok{70e{-}6}\NormalTok{,}
\NormalTok{    gate\_time\_1q: }\BuiltInTok{float} \OperatorTok{=} \FloatTok{50e{-}9}\NormalTok{,}
\NormalTok{    gate\_time\_2q: }\BuiltInTok{float} \OperatorTok{=} \FloatTok{300e{-}9}\NormalTok{,}
\NormalTok{) }\OperatorTok{{-}\textgreater{}}\NormalTok{ qml.Device:}
    \CommentTok{"""Return a PennyLane default.mixed device for mixed{-}state simulation.}

\CommentTok{    The device itself is noise{-}agnostic; noise channels (depolarising,}
\CommentTok{    thermal relaxation) are inserted explicitly inside the QNode via}
\CommentTok{    channel operations.}

\CommentTok{    Returns:}
\CommentTok{        PennyLane default.mixed device with n\_qubits wires.}
\CommentTok{    """}
\NormalTok{    ...}


\KeywordTok{def}\NormalTok{ noisy\_qaoa\_qnode(}
\NormalTok{    cost\_hamiltonian: qml.Hamiltonian,}
\NormalTok{    n\_qubits: }\BuiltInTok{int}\NormalTok{,}
\NormalTok{    p: }\BuiltInTok{int}\NormalTok{,}
\NormalTok{    noise\_params: }\BuiltInTok{dict}\NormalTok{,}
\NormalTok{) }\OperatorTok{{-}\textgreater{}}\NormalTok{ qml.QNode:}
    \CommentTok{"""Construct a QNode on default.mixed that interleaves noise channels.}

\CommentTok{    After each single{-}qubit gate: DepolarizingChannel + ThermalRelaxationError.}
\CommentTok{    After each two{-}qubit IsingZZ: both qubits receive depolarising noise}
\CommentTok{    at \textasciitilde{}10x the single{-}qubit rate (reflecting typical hardware) and}
\CommentTok{    thermal relaxation scaled by the two{-}qubit gate duration.}

\CommentTok{    Args:}
\CommentTok{        noise\_params: Dict with keys \textquotesingle{}depol\_rate\textquotesingle{}, \textquotesingle{}t1\textquotesingle{}, \textquotesingle{}t2\textquotesingle{},}
\CommentTok{                      \textquotesingle{}gate\_time\_1q\textquotesingle{}, \textquotesingle{}gate\_time\_2q\textquotesingle{}.}

\CommentTok{    Returns:}
\CommentTok{        QNode producing \textless{}H\_C\textgreater{} under gate{-}dependent incoherent noise.}
\CommentTok{    """}
\NormalTok{    ...}


\KeywordTok{def}\NormalTok{ lindblad\_simulation(}
\NormalTok{    initial\_state: qutip.Qobj,}
\NormalTok{    hamiltonian: qutip.Qobj,}
\NormalTok{    collapse\_ops: }\BuiltInTok{list}\NormalTok{[qutip.Qobj],}
\NormalTok{    tlist: np.ndarray,}
\NormalTok{    observables: }\BuiltInTok{dict}\NormalTok{[}\BuiltInTok{str}\NormalTok{, qutip.Qobj],}
\NormalTok{) }\OperatorTok{{-}\textgreater{}} \BuiltInTok{dict}\NormalTok{[}\BuiltInTok{str}\NormalTok{, np.ndarray]:}
    \CommentTok{"""Run a full Lindblad master equation simulation via qutip.mesolve.}

\CommentTok{    Evolves the density matrix under:}
\CommentTok{        drho/dt = {-}i[H, rho] + sum\_k (L\_k rho L\_k\^{}dag {-} 1/2 \{L\_k\^{}dag L\_k, rho\})}

\CommentTok{    Collapse operators encode amplitude damping (L = sqrt(gamma\_1) sigma\_{-})}
\CommentTok{    and pure dephasing (L = sqrt(gamma\_phi) Z/2).}

\CommentTok{    Returns:}
\CommentTok{        Dict mapping observable labels to time{-}series arrays of}
\CommentTok{        shape (len(tlist),).}
\CommentTok{    """}
\NormalTok{    ...}


\KeywordTok{def}\NormalTok{ compare\_noise\_channels(}
\NormalTok{    n\_qubits: }\BuiltInTok{int}\NormalTok{,}
\NormalTok{    circuit\_params: np.ndarray,}
\NormalTok{    noise\_params: }\BuiltInTok{dict}\NormalTok{,}
\NormalTok{) }\OperatorTok{{-}\textgreater{}} \BuiltInTok{dict}\NormalTok{[}\BuiltInTok{str}\NormalTok{, }\BuiltInTok{float}\NormalTok{]:}
    \CommentTok{"""Compare PennyLane channel model vs QuTiP Lindblad for the same circuit.}

\CommentTok{    Runs the same QAOA circuit under both models and compares:}
\CommentTok{    {-} Energy expectation values}
\CommentTok{    {-} Inter{-}model state fidelity}
\CommentTok{    {-} Output{-}state purity Tr(rho\^{}2)}

\CommentTok{    Returns:}
\CommentTok{        Dict with keys \textquotesingle{}energy\_pennylane\textquotesingle{}, \textquotesingle{}energy\_qutip\textquotesingle{}, \textquotesingle{}fidelity\textquotesingle{},}
\CommentTok{        \textquotesingle{}purity\_pl\textquotesingle{}, \textquotesingle{}purity\_qt\textquotesingle{}.}
\CommentTok{    """}
\NormalTok{    ...}
\end{Highlighting}
\end{Shaded}

\emph{(Full implementation: \texttt{report1\_scripts/noise\_model.py})}

\subsubsection{\texorpdfstring{6.
\texttt{hardware\_runner.py}}{6. hardware\_runner.py}}\label{hardware_runner.py}

\textbf{Purpose:} Submit QAOA circuits to IBM Quantum or IQM via
PennyLane device plugins.

\begin{Shaded}
\begin{Highlighting}[]
\ImportTok{from}\NormalTok{ \_\_future\_\_ }\ImportTok{import}\NormalTok{ annotations}
\ImportTok{from}\NormalTok{ dataclasses }\ImportTok{import}\NormalTok{ dataclass}
\ImportTok{import}\NormalTok{ numpy }\ImportTok{as}\NormalTok{ np}
\ImportTok{import}\NormalTok{ pennylane }\ImportTok{as}\NormalTok{ qml}


\AttributeTok{@dataclass}
\KeywordTok{class}\NormalTok{ HardwareResult:}
    \CommentTok{"""Container for hardware execution results.}

\CommentTok{    Attributes:}
\CommentTok{        expectation\_value: Measured \textless{}H\_C\textgreater{}.}
\CommentTok{        bitstring\_counts: Raw bitstring counts dict.}
\CommentTok{        shots: Number of shots used.}
\CommentTok{        backend\_name: Name of the quantum backend.}
\CommentTok{        wall\_time: Wall{-}clock execution time in seconds.}
\CommentTok{    """}
\NormalTok{    expectation\_value: }\BuiltInTok{float}
\NormalTok{    bitstring\_counts: }\BuiltInTok{dict}\NormalTok{[}\BuiltInTok{str}\NormalTok{, }\BuiltInTok{int}\NormalTok{]}
\NormalTok{    shots: }\BuiltInTok{int}
\NormalTok{    backend\_name: }\BuiltInTok{str}
\NormalTok{    wall\_time: }\BuiltInTok{float}


\KeywordTok{def}\NormalTok{ create\_ibm\_device(}
\NormalTok{    n\_qubits: }\BuiltInTok{int}\NormalTok{,}
\NormalTok{    backend\_name: }\BuiltInTok{str} \OperatorTok{=} \StringTok{"ibm\_brisbane"}\NormalTok{,}
\NormalTok{    shots: }\BuiltInTok{int} \OperatorTok{=} \DecValTok{4096}\NormalTok{,}
\NormalTok{    ibm\_token: }\BuiltInTok{str} \OperatorTok{|} \VariableTok{None} \OperatorTok{=} \VariableTok{None}\NormalTok{,}
\NormalTok{) }\OperatorTok{{-}\textgreater{}}\NormalTok{ qml.Device:}
    \CommentTok{"""Return a PennyLane device targeting IBM Quantum hardware.}

\CommentTok{    Uses the PennyLane{-}Qiskit plugin (qml.device(\textquotesingle{}qiskit.ibmq\textquotesingle{}, ...)).}
\CommentTok{    The API token is read from the IBMQ\_TOKEN environment variable if}
\CommentTok{    not supplied explicitly. PennyLane handles transpilation internally.}

\CommentTok{    Returns:}
\CommentTok{        PennyLane device bound to an IBM Quantum backend.}
\CommentTok{    """}
\NormalTok{    ...}


\KeywordTok{def}\NormalTok{ create\_iqm\_device(}
\NormalTok{    n\_qubits: }\BuiltInTok{int}\NormalTok{,}
\NormalTok{    backend\_url: }\BuiltInTok{str}\NormalTok{,}
\NormalTok{    shots: }\BuiltInTok{int} \OperatorTok{=} \DecValTok{4096}\NormalTok{,}
\NormalTok{) }\OperatorTok{{-}\textgreater{}}\NormalTok{ qml.Device:}
    \CommentTok{"""Return a PennyLane device targeting IQM hardware.}

\CommentTok{    Uses the PennyLane{-}IQM plugin.}

\CommentTok{    Returns:}
\CommentTok{        PennyLane device bound to an IQM backend.}
\CommentTok{    """}
\NormalTok{    ...}


\KeywordTok{def}\NormalTok{ execute\_on\_hardware(}
\NormalTok{    qnode: qml.QNode,}
\NormalTok{    params: np.ndarray,}
\NormalTok{    cost\_hamiltonian: qml.Hamiltonian,}
\NormalTok{) }\OperatorTok{{-}\textgreater{}}\NormalTok{ HardwareResult:}
    \CommentTok{"""Run the QAOA circuit on hardware and measure \textless{}H\_C\textgreater{}.}

\CommentTok{    Two measurement strategies:}
\CommentTok{        1. Direct: qml.expval(H\_C) {-}{-} PennyLane decomposes the}
\CommentTok{           Hamiltonian into commuting Pauli groups.}
\CommentTok{        2. Sample{-}based: qml.sample() returns raw bitstrings of}
\CommentTok{           shape (S, N). Energy per sample computed classically via}
\CommentTok{           C(z) = sum\_\{(i,j)\} w\_ij/2 * (1 {-} s\_i * s\_j) where s\_i = 1 {-} 2*z\_i.}

\CommentTok{    Returns:}
\CommentTok{        HardwareResult with expectation value, counts, and metadata.}
\CommentTok{    """}
\NormalTok{    ...}


\KeywordTok{def}\NormalTok{ select\_optimal\_qubits(}
\NormalTok{    backend\_name: }\BuiltInTok{str}\NormalTok{,}
\NormalTok{    n\_qubits: }\BuiltInTok{int}\NormalTok{,}
\NormalTok{    ibm\_token: }\BuiltInTok{str} \OperatorTok{|} \VariableTok{None} \OperatorTok{=} \VariableTok{None}\NormalTok{,}
\NormalTok{) }\OperatorTok{{-}\textgreater{}} \BuiltInTok{list}\NormalTok{[}\BuiltInTok{int}\NormalTok{]:}
    \CommentTok{"""Select the best{-}performing connected subgraph of physical qubits.}

\CommentTok{    Loads backend calibration data (T1, T2, CNOT error rates) and builds}
\CommentTok{    a weighted graph over physical qubits (edges weighted by 1/eps\_CX).}
\CommentTok{    Searches for the size{-}N connected subgraph maximising}
\CommentTok{        sum\_edges 1/eps\_CX + alpha * sum\_nodes T1*T2}
\CommentTok{    via BFS enumeration.}

\CommentTok{    Returns:}
\CommentTok{        List of physical qubit indices for the initial\_layout.}
\CommentTok{    """}
\NormalTok{    ...}
\end{Highlighting}
\end{Shaded}

\emph{(Full implementation:
\texttt{report1\_scripts/hardware\_runner.py})}

\subsubsection{\texorpdfstring{7.
\texttt{error\_mitigation.py}}{7. error\_mitigation.py}}\label{error_mitigation.py}

\textbf{Purpose:} Zero-noise extrapolation (ZNE) and supporting error
mitigation techniques.

\begin{Shaded}
\begin{Highlighting}[]
\ImportTok{from}\NormalTok{ \_\_future\_\_ }\ImportTok{import}\NormalTok{ annotations}
\ImportTok{from}\NormalTok{ typing }\ImportTok{import}\NormalTok{ Callable, Literal}
\ImportTok{import}\NormalTok{ numpy }\ImportTok{as}\NormalTok{ np}
\ImportTok{import}\NormalTok{ pennylane }\ImportTok{as}\NormalTok{ qml}


\KeywordTok{def}\NormalTok{ zne\_extrapolate(}
\NormalTok{    qnode: qml.QNode,}
\NormalTok{    params: np.ndarray,}
\NormalTok{    scale\_factors: }\BuiltInTok{list}\NormalTok{[}\BuiltInTok{float}\NormalTok{] }\OperatorTok{=}\NormalTok{ [}\FloatTok{1.0}\NormalTok{, }\FloatTok{3.0}\NormalTok{, }\FloatTok{5.0}\NormalTok{],}
\NormalTok{    extrapolation: Literal[}\StringTok{"linear"}\NormalTok{, }\StringTok{"richardson"}\NormalTok{] }\OperatorTok{=} \StringTok{"richardson"}\NormalTok{,}
\NormalTok{) }\OperatorTok{{-}\textgreater{}} \BuiltInTok{tuple}\NormalTok{[}\BuiltInTok{float}\NormalTok{, }\BuiltInTok{dict}\NormalTok{]:}
    \CommentTok{"""Apply zero{-}noise extrapolation to a hardware or noisy{-}simulator QNode.}

\CommentTok{    For each scale factor lambda, amplifies circuit noise via global}
\CommentTok{    unitary folding and measures \textless{}H\_C\textgreater{}\_lambda. Extrapolates the}
\CommentTok{    (lambda, E(lambda)) data to lambda=0.}

\CommentTok{    Returns:}
\CommentTok{        Tuple of (extrapolated\_value, details\_dict) where details\_dict}
\CommentTok{        contains \textquotesingle{}noisy\_values\textquotesingle{}, \textquotesingle{}scale\_factors\textquotesingle{}, and \textquotesingle{}fit\_coefficients\textquotesingle{}.}
\CommentTok{    """}
\NormalTok{    ...}


\KeywordTok{def}\NormalTok{ global\_fold(}
\NormalTok{    circuit\_fn: Callable,}
\NormalTok{    scale\_factor: }\BuiltInTok{float}\NormalTok{,}
\NormalTok{) }\OperatorTok{{-}\textgreater{}}\NormalTok{ Callable:}
    \CommentTok{"""Produce a noise{-}amplified circuit by global unitary folding.}

\CommentTok{    For scale factor lambda, replaces circuit U with U(U\^{}dag U)\^{}n where}
\CommentTok{    n = floor((lambda{-}1)/2). For non{-}integer scale factors, partial}
\CommentTok{    folding applies to the last frac * N\_gates gates where}
\CommentTok{    frac = (lambda{-}1)/2 {-} n.}

\CommentTok{    Returns:}
\CommentTok{        A callable that applies the folded circuit.}
\CommentTok{    """}
\NormalTok{    ...}


\KeywordTok{def}\NormalTok{ richardson\_extrapolation(}
\NormalTok{    scale\_factors: np.ndarray,}
\NormalTok{    noisy\_values: np.ndarray,}
\NormalTok{) }\OperatorTok{{-}\textgreater{}} \BuiltInTok{float}\NormalTok{:}
    \CommentTok{"""Fit a degree{-}(M{-}1) polynomial through M data points via Vandermonde.}

\CommentTok{    Solves V @ a = E for the coefficient vector a, where V is the}
\CommentTok{    Vandermonde matrix of the scale factors. Returns a\_0 = E(0).}

\CommentTok{    Returns:}
\CommentTok{        Extrapolated zero{-}noise energy value.}
\CommentTok{    """}
\NormalTok{    ...}


\KeywordTok{def}\NormalTok{ linear\_extrapolation(}
\NormalTok{    scale\_factors: np.ndarray,}
\NormalTok{    noisy\_values: np.ndarray,}
\NormalTok{) }\OperatorTok{{-}\textgreater{}} \BuiltInTok{float}\NormalTok{:}
    \CommentTok{"""Fit E(lambda) = a + b*lambda via least{-}squares.}

\CommentTok{    Returns:}
\CommentTok{        a = E(0), the extrapolated zero{-}noise energy value.}
\CommentTok{    """}
\NormalTok{    ...}
\end{Highlighting}
\end{Shaded}

\emph{(Full implementation:
\texttt{report1\_scripts/error\_mitigation.py})}

\subsubsection{\texorpdfstring{8.
\texttt{renderer.py}}{8. renderer.py}}\label{renderer.py}

\textbf{Purpose:} OpenGL 3D surface rendering of the
\((\beta, \gamma, \langle H_C \rangle)\) landscape.

\begin{Shaded}
\begin{Highlighting}[]
\ImportTok{from}\NormalTok{ \_\_future\_\_ }\ImportTok{import}\NormalTok{ annotations}
\ImportTok{import}\NormalTok{ numpy }\ImportTok{as}\NormalTok{ np}
\ImportTok{import}\NormalTok{ glfw}
\ImportTok{from}\NormalTok{ OpenGL.GL }\ImportTok{import} \OperatorTok{*}


\KeywordTok{def}\NormalTok{ init\_opengl\_context(}
\NormalTok{    width: }\BuiltInTok{int} \OperatorTok{=} \DecValTok{1280}\NormalTok{,}
\NormalTok{    height: }\BuiltInTok{int} \OperatorTok{=} \DecValTok{720}\NormalTok{,}
\NormalTok{    title: }\BuiltInTok{str} \OperatorTok{=} \StringTok{"QAOA Energy Landscape"}\NormalTok{,}
\NormalTok{) }\OperatorTok{{-}\textgreater{}}\NormalTok{ glfw.\_GLFWwindow:}
    \CommentTok{"""Initialise a GLFW window with OpenGL 3.3 core profile.}

\CommentTok{    Enables depth testing, 4x MSAA, and alpha blending.}

\CommentTok{    Returns:}
\CommentTok{        The GLFW window handle.}
\CommentTok{    """}
\NormalTok{    ...}


\KeywordTok{def}\NormalTok{ compile\_shader\_program(}
\NormalTok{    vertex\_source: }\BuiltInTok{str}\NormalTok{,}
\NormalTok{    fragment\_source: }\BuiltInTok{str}\NormalTok{,}
\NormalTok{) }\OperatorTok{{-}\textgreater{}} \BuiltInTok{int}\NormalTok{:}
    \CommentTok{"""Compile and link a vertex{-}fragment shader pair.}

\CommentTok{    Returns:}
\CommentTok{        OpenGL shader program ID.}
\CommentTok{    """}
\NormalTok{    ...}


\KeywordTok{def}\NormalTok{ create\_surface\_vbo(}
\NormalTok{    gamma\_vals: np.ndarray,}
\NormalTok{    beta\_vals: np.ndarray,}
\NormalTok{    energies: np.ndarray,}
\NormalTok{) }\OperatorTok{{-}\textgreater{}} \BuiltInTok{tuple}\NormalTok{[}\BuiltInTok{int}\NormalTok{, }\BuiltInTok{int}\NormalTok{, }\BuiltInTok{int}\NormalTok{, }\BuiltInTok{int}\NormalTok{]:}
    \CommentTok{"""Build a triangle{-}mesh VAO/VBO/EBO from the (N\_gamma x N\_beta) grid.}

\CommentTok{    Each vertex stores 7 float32 values (28 bytes): position}
\CommentTok{    (gamma, beta, \textless{}H\_C\textgreater{}), surface normal (n\_x, n\_y, n\_z) from}
\CommentTok{    finite{-}difference gradients, and normalised energy in [0,1].}
\CommentTok{    Two triangles per grid cell. Uses GL\_DYNAMIC\_DRAW for streaming.}

\CommentTok{    Returns:}
\CommentTok{        Tuple of (vao\_id, vbo\_id, ebo\_id, n\_indices).}
\CommentTok{    """}
\NormalTok{    ...}


\KeywordTok{def}\NormalTok{ update\_surface\_vbo(}
\NormalTok{    vbo\_id: }\BuiltInTok{int}\NormalTok{,}
\NormalTok{    gamma\_vals: np.ndarray,}
\NormalTok{    beta\_vals: np.ndarray,}
\NormalTok{    energies: np.ndarray,}
\NormalTok{) }\OperatorTok{{-}\textgreater{}} \VariableTok{None}\NormalTok{:}
    \CommentTok{"""Overwrite existing VBO vertex data in{-}place via glBufferSubData."""}
\NormalTok{    ...}


\KeywordTok{def}\NormalTok{ draw\_optimizer\_trajectory(}
\NormalTok{    trajectory: }\BuiltInTok{list}\NormalTok{[}\BuiltInTok{tuple}\NormalTok{[}\BuiltInTok{float}\NormalTok{, }\BuiltInTok{float}\NormalTok{, }\BuiltInTok{float}\NormalTok{]],}
\NormalTok{    shader\_program: }\BuiltInTok{int}\NormalTok{,}
\NormalTok{) }\OperatorTok{{-}\textgreater{}} \VariableTok{None}\NormalTok{:}
    \CommentTok{"""Render the optimiser\textquotesingle{}s path as a GL\_LINE\_STRIP above the surface.}

\CommentTok{    Each vertex carries a normalised progress attribute t in [0,1]}
\CommentTok{    for blue{-}to{-}red colour interpolation. Slight positive z{-}offset}
\CommentTok{    (\textasciitilde{}0.01) prevents z{-}fighting. Uses GL\_STREAM\_DRAW.}
\CommentTok{    """}
\NormalTok{    ...}


\KeywordTok{def}\NormalTok{ render\_side\_panel(}
\NormalTok{    p: }\BuiltInTok{int}\NormalTok{,}
\NormalTok{    current\_energy: }\BuiltInTok{float}\NormalTok{,}
\NormalTok{    approximation\_ratio: }\BuiltInTok{float}\NormalTok{,}
\NormalTok{    gradient\_norm: }\BuiltInTok{float}\NormalTok{,}
\NormalTok{    iteration: }\BuiltInTok{int}\NormalTok{,}
\NormalTok{) }\OperatorTok{{-}\textgreater{}} \VariableTok{None}\NormalTok{:}
    \CommentTok{"""Draw a text overlay with live optimisation statistics."""}
\NormalTok{    ...}


\KeywordTok{class}\NormalTok{ ArcballCamera:}
    \CommentTok{"""Interactive 3D camera with arcball rotation, zoom, and pan.}

\CommentTok{    Mouse drag rotates around a focus point (yaw/pitch angles),}
\CommentTok{    scroll wheel zooms by adjusting camera{-}target distance,}
\CommentTok{    middle{-}click pans. Pitch clamped to +/{-}89 degrees.}
\CommentTok{    """}

    \KeywordTok{def} \FunctionTok{\_\_init\_\_}\NormalTok{(}
        \VariableTok{self}\NormalTok{,}
\NormalTok{        target: np.ndarray }\OperatorTok{=}\NormalTok{ np.zeros(}\DecValTok{3}\NormalTok{),}
\NormalTok{        distance: }\BuiltInTok{float} \OperatorTok{=} \FloatTok{5.0}\NormalTok{,}
\NormalTok{        yaw: }\BuiltInTok{float} \OperatorTok{=} \OperatorTok{{-}}\FloatTok{90.0}\NormalTok{,}
\NormalTok{        pitch: }\BuiltInTok{float} \OperatorTok{=} \FloatTok{30.0}\NormalTok{,}
\NormalTok{    ) }\OperatorTok{{-}\textgreater{}} \VariableTok{None}\NormalTok{:}
\NormalTok{        ...}

    \KeywordTok{def}\NormalTok{ get\_view\_matrix(}\VariableTok{self}\NormalTok{) }\OperatorTok{{-}\textgreater{}}\NormalTok{ np.ndarray:}
        \CommentTok{"""Return a 4x4 look{-}at view matrix (float32)."""}
\NormalTok{        ...}

    \KeywordTok{def}\NormalTok{ get\_projection\_matrix(}
        \VariableTok{self}\NormalTok{,}
\NormalTok{        aspect\_ratio: }\BuiltInTok{float}\NormalTok{,}
\NormalTok{        fov: }\BuiltInTok{float} \OperatorTok{=} \FloatTok{45.0}\NormalTok{,}
\NormalTok{        near: }\BuiltInTok{float} \OperatorTok{=} \FloatTok{0.1}\NormalTok{,}
\NormalTok{        far: }\BuiltInTok{float} \OperatorTok{=} \FloatTok{100.0}\NormalTok{,}
\NormalTok{    ) }\OperatorTok{{-}\textgreater{}}\NormalTok{ np.ndarray:}
        \CommentTok{"""Return a 4x4 perspective projection matrix (float32)."""}
\NormalTok{        ...}
\end{Highlighting}
\end{Shaded}

\emph{(Full implementation: \texttt{report1\_scripts/renderer.py})}

\subsubsection{\texorpdfstring{9.
\texttt{main.py}}{9. main.py}}\label{main.py}

\textbf{Purpose:} CLI entry point and workflow orchestration.

\begin{Shaded}
\begin{Highlighting}[]
\ImportTok{from}\NormalTok{ \_\_future\_\_ }\ImportTok{import}\NormalTok{ annotations}
\ImportTok{import}\NormalTok{ argparse}
\ImportTok{from}\NormalTok{ pathlib }\ImportTok{import}\NormalTok{ Path}


\KeywordTok{def}\NormalTok{ parse\_args() }\OperatorTok{{-}\textgreater{}}\NormalTok{ argparse.Namespace:}
    \CommentTok{"""Define and parse CLI arguments.}

\CommentTok{    Key flags:}
\CommentTok{        {-}{-}graph{-}type    Graph family (cycle, complete, erdos\_renyi, etc.)}
\CommentTok{        {-}{-}n{-}nodes       Number of graph nodes}
\CommentTok{        {-}{-}p             QAOA depth (number of layers)}
\CommentTok{        {-}{-}optimizer     Optimiser backend (COBYLA / L{-}BFGS{-}B / SPSA / Adam)}
\CommentTok{        {-}{-}grid{-}size     Landscape sampling resolution}
\CommentTok{        {-}{-}noise         Enable noisy simulation (default.mixed)}
\CommentTok{        {-}{-}hardware      Run on a real quantum processor}
\CommentTok{        {-}{-}backend       Backend name (e.g. ibm\_brisbane)}
\CommentTok{        {-}{-}zne           Enable zero{-}noise extrapolation}
\CommentTok{        {-}{-}render        Launch OpenGL visualisation}
\CommentTok{        {-}{-}config        Path to a YAML configuration file}
\CommentTok{        {-}{-}seed          Random seed for reproducibility}

\CommentTok{    Returns:}
\CommentTok{        Parsed arguments namespace.}
\CommentTok{    """}
\NormalTok{    ...}


\KeywordTok{def}\NormalTok{ run\_pipeline(args: argparse.Namespace) }\OperatorTok{{-}\textgreater{}} \VariableTok{None}\NormalTok{:}
    \CommentTok{"""Execute the full QAOA exploration pipeline.}

\CommentTok{    Steps:}
\CommentTok{        1. Graph generation {-}{-} create\_graph + build\_cost\_hamiltonian;}
\CommentTok{           compute C\_max via brute force (N \textless{}= 20).}
\CommentTok{        2. Circuit construction {-}{-} build QNode on appropriate backend}
\CommentTok{           (noiseless simulator, noisy default.mixed, or hardware).}
\CommentTok{        3. Landscape sampling {-}{-} for p=1 with rendering, sample the}
\CommentTok{           2D (gamma, beta) energy grid.}
\CommentTok{        4. Classical optimisation {-}{-} run optimize\_qaoa, recording the}
\CommentTok{           trajectory for visualisation.}
\CommentTok{        5. Error mitigation {-}{-} if ZNE requested, apply zne\_extrapolate}
\CommentTok{           at optimal parameters.}
\CommentTok{        6. Rendering {-}{-} if enabled, launch OpenGL window with surface,}
\CommentTok{           trajectory overlay, and status panel.}
\CommentTok{    """}
\NormalTok{    ...}
\end{Highlighting}
\end{Shaded}

\emph{(Full implementation: \texttt{report1\_scripts/main.py})}

\begin{center}\rule{0.5\linewidth}{0.5pt}\end{center}

\section{Part 3 --- Implementation Guidelines (Instructional
Pseudocode)}\label{part-3-implementation-guidelines-instructional-pseudocode}

\subsection{\texorpdfstring{Module 1:
\texttt{graph\_generator.py}}{Module 1: graph\_generator.py}}\label{module-1-graph_generator.py}

\begin{Shaded}
\begin{Highlighting}[]
\CommentTok{\# {-}{-}{-} graph\_generator.py pseudocode {-}{-}{-}}

\KeywordTok{def}\NormalTok{ create\_graph(graph\_type, n\_nodes, seed, }\OperatorTok{**}\NormalTok{kwargs):}
    \CommentTok{\# Dispatch on graph\_type string}
    \ControlFlowTok{if}\NormalTok{ graph\_type }\OperatorTok{==} \StringTok{"cycle"}\NormalTok{:}
\NormalTok{        G }\OperatorTok{=}\NormalTok{ nx.cycle\_graph(n\_nodes)}
    \ControlFlowTok{elif}\NormalTok{ graph\_type }\OperatorTok{==} \StringTok{"complete"}\NormalTok{:}
\NormalTok{        G }\OperatorTok{=}\NormalTok{ nx.complete\_graph(n\_nodes)}
    \ControlFlowTok{elif}\NormalTok{ graph\_type }\OperatorTok{==} \StringTok{"erdos\_renyi"}\NormalTok{:}
\NormalTok{        G }\OperatorTok{=}\NormalTok{ nx.erdos\_renyi\_graph(n\_nodes, kwargs[}\StringTok{"p"}\NormalTok{], seed}\OperatorTok{=}\NormalTok{seed)}
    \ControlFlowTok{elif}\NormalTok{ graph\_type }\OperatorTok{==} \StringTok{"random\_regular"}\NormalTok{:}
\NormalTok{        G }\OperatorTok{=}\NormalTok{ nx.random\_regular\_graph(kwargs[}\StringTok{"d"}\NormalTok{], n\_nodes, seed}\OperatorTok{=}\NormalTok{seed)}
    \ControlFlowTok{elif}\NormalTok{ graph\_type }\OperatorTok{==} \StringTok{"grid\_2d"}\NormalTok{:}
\NormalTok{        rows, cols }\OperatorTok{=}\NormalTok{ kwargs.get(}\StringTok{"rows"}\NormalTok{, n\_nodes), kwargs.get(}\StringTok{"cols"}\NormalTok{, n\_nodes)}
\NormalTok{        G }\OperatorTok{=}\NormalTok{ nx.grid\_2d\_graph(rows, cols)}
\NormalTok{        G }\OperatorTok{=}\NormalTok{ nx.convert\_node\_labels\_to\_integers(G)  }\CommentTok{\# relabel (i,j) {-}\textgreater{} int}
    \ControlFlowTok{elif}\NormalTok{ graph\_type }\OperatorTok{==} \StringTok{"petersen"}\NormalTok{:}
\NormalTok{        G }\OperatorTok{=}\NormalTok{ nx.petersen\_graph()}
    \ControlFlowTok{elif}\NormalTok{ graph\_type }\OperatorTok{==} \StringTok{"custom"}\NormalTok{:}
\NormalTok{        G }\OperatorTok{=}\NormalTok{ nx.Graph()}
\NormalTok{        G.add\_edges\_from(kwargs[}\StringTok{"edges"}\NormalTok{])}
    \CommentTok{\# ... other types ...}

    \CommentTok{\# Assign unit weights to edges that don\textquotesingle{}t have a weight}
    \ControlFlowTok{for}\NormalTok{ u, v }\KeywordTok{in}\NormalTok{ G.edges():}
        \ControlFlowTok{if} \StringTok{"weight"} \KeywordTok{not} \KeywordTok{in}\NormalTok{ G[u][v]:}
\NormalTok{            G[u][v][}\StringTok{"weight"}\NormalTok{] }\OperatorTok{=} \FloatTok{1.0}
    \ControlFlowTok{return}\NormalTok{ G}


\KeywordTok{def}\NormalTok{ adjacency\_to\_ising(graph):}
\NormalTok{    coeffs }\OperatorTok{=}\NormalTok{ []                         }\CommentTok{\# list of float}
\NormalTok{    obs }\OperatorTok{=}\NormalTok{ []                            }\CommentTok{\# list of (int, int) qubit pairs}
\NormalTok{    offset }\OperatorTok{=} \FloatTok{0.0}
    \ControlFlowTok{for}\NormalTok{ i, j, w }\KeywordTok{in}\NormalTok{ graph.edges(data}\OperatorTok{=}\StringTok{"weight"}\NormalTok{):}
        \CommentTok{\# MaxCut: H\_C = sum w\_ij/2 (I {-} Z\_i Z\_j)}
\NormalTok{        coeffs.append(}\OperatorTok{{-}}\NormalTok{w }\OperatorTok{/} \DecValTok{2}\NormalTok{)}
\NormalTok{        obs.append((i, j))              }\CommentTok{\# Pauli Z\_i Z\_j pair}
\NormalTok{        offset }\OperatorTok{+=}\NormalTok{ w }\OperatorTok{/} \DecValTok{2}                 \CommentTok{\# identity contribution}
    \ControlFlowTok{return}\NormalTok{ coeffs, obs, offset}


\KeywordTok{def}\NormalTok{ build\_cost\_hamiltonian(graph):}
\NormalTok{    coeffs, obs, offset }\OperatorTok{=}\NormalTok{ adjacency\_to\_ising(graph)}
    \CommentTok{\# Build PennyLane Hamiltonian}
\NormalTok{    all\_coeffs }\OperatorTok{=}\NormalTok{ coeffs }\OperatorTok{+}\NormalTok{ [offset]}
\NormalTok{    all\_obs }\OperatorTok{=}\NormalTok{ [qml.PauliZ(i) }\OperatorTok{@}\NormalTok{ qml.PauliZ(j) }\ControlFlowTok{for}\NormalTok{ i, j }\KeywordTok{in}\NormalTok{ obs] }\OperatorTok{+}\NormalTok{ [qml.Identity(}\DecValTok{0}\NormalTok{)]}
    \ControlFlowTok{return}\NormalTok{ qml.Hamiltonian(all\_coeffs, all\_obs)}
\end{Highlighting}
\end{Shaded}

\emph{(Full implementation:
\texttt{report1\_scripts/graph\_generator.py})}

\subsection{\texorpdfstring{Module 2:
\texttt{qaoa\_circuit.py}}{Module 2: qaoa\_circuit.py}}\label{module-2-qaoa_circuit.py}

\begin{Shaded}
\begin{Highlighting}[]
\CommentTok{\# {-}{-}{-} qaoa\_circuit.py pseudocode {-}{-}{-}}

\KeywordTok{def}\NormalTok{ cost\_layer(gamma, graph\_edges):}
    \CommentTok{\# graph\_edges: list of (i, j, w\_ij)}
    \ControlFlowTok{for}\NormalTok{ i, j, w }\KeywordTok{in}\NormalTok{ graph\_edges:}
        \CommentTok{\# IsingZZ implements exp({-}i phi Z⊗Z / 2), so pass phi = gamma * w}
\NormalTok{        qml.IsingZZ(gamma }\OperatorTok{*}\NormalTok{ w, wires}\OperatorTok{=}\NormalTok{[i, j])}


\KeywordTok{def}\NormalTok{ mixer\_layer(beta, n\_qubits):}
    \ControlFlowTok{for}\NormalTok{ q }\KeywordTok{in} \BuiltInTok{range}\NormalTok{(n\_qubits):}
        \CommentTok{\# RX(theta) = exp({-}i theta X / 2), so theta = 2*beta for exp({-}i beta X)}
\NormalTok{        qml.RX(}\DecValTok{2} \OperatorTok{*}\NormalTok{ beta, wires}\OperatorTok{=}\NormalTok{q)}


\KeywordTok{def}\NormalTok{ extract\_edges\_from\_hamiltonian(cost\_hamiltonian):}
    \CommentTok{"""Recover weighted edge list from qml.Hamiltonian."""}
\NormalTok{    edges }\OperatorTok{=}\NormalTok{ []}
    \ControlFlowTok{for}\NormalTok{ coeff, obs }\KeywordTok{in} \BuiltInTok{zip}\NormalTok{(cost\_hamiltonian.coeffs, cost\_hamiltonian.ops):}
        \ControlFlowTok{if} \BuiltInTok{isinstance}\NormalTok{(obs, qml.Identity):}
            \ControlFlowTok{continue}  \CommentTok{\# skip the offset term}
        \CommentTok{\# obs is PauliZ(i) @ PauliZ(j); coeff is {-}w/2}
\NormalTok{        wires }\OperatorTok{=}\NormalTok{ obs.wires.tolist()       }\CommentTok{\# [i, j]}
\NormalTok{        w }\OperatorTok{=} \OperatorTok{{-}}\DecValTok{2} \OperatorTok{*}\NormalTok{ coeff                   }\CommentTok{\# recover original weight}
\NormalTok{        edges.append((wires[}\DecValTok{0}\NormalTok{], wires[}\DecValTok{1}\NormalTok{], w))}
    \ControlFlowTok{return}\NormalTok{ edges}


\KeywordTok{def}\NormalTok{ qaoa\_ansatz(params, cost\_hamiltonian, n\_qubits, p):}
    \CommentTok{\# params shape: (2p,) {-}{-} [gamma\_1, ..., gamma\_p, beta\_1, ..., beta\_p]}
\NormalTok{    gammas }\OperatorTok{=}\NormalTok{ params[:p]                  }\CommentTok{\# shape (p,)}
\NormalTok{    betas }\OperatorTok{=}\NormalTok{ params[p:]                   }\CommentTok{\# shape (p,)}
\NormalTok{    edges }\OperatorTok{=}\NormalTok{ extract\_edges\_from\_hamiltonian(cost\_hamiltonian)}

    \CommentTok{\# 1. Prepare |+\textgreater{}\^{}N}
    \ControlFlowTok{for}\NormalTok{ q }\KeywordTok{in} \BuiltInTok{range}\NormalTok{(n\_qubits):}
\NormalTok{        qml.Hadamard(wires}\OperatorTok{=}\NormalTok{q)}

    \CommentTok{\# 2. Apply p layers of cost + mixer}
    \ControlFlowTok{for}\NormalTok{ k }\KeywordTok{in} \BuiltInTok{range}\NormalTok{(p):}
\NormalTok{        cost\_layer(gammas[k], edges)     }\CommentTok{\# U\_C(gamma\_k)}
\NormalTok{        mixer\_layer(betas[k], n\_qubits)  }\CommentTok{\# U\_B(beta\_k)}


\KeywordTok{def}\NormalTok{ build\_qaoa\_circuit(cost\_hamiltonian, n\_qubits, p, device):}
    \AttributeTok{@qml.qnode}\NormalTok{(device, diff\_method}\OperatorTok{=}\StringTok{"best"}\NormalTok{)}
    \KeywordTok{def}\NormalTok{ circuit(params):}
        \CommentTok{\# params: np.ndarray of shape (2p,)}
\NormalTok{        qaoa\_ansatz(params, cost\_hamiltonian, n\_qubits, p)}
        \ControlFlowTok{return}\NormalTok{ qml.expval(cost\_hamiltonian)}
    \ControlFlowTok{return}\NormalTok{ circuit}
\end{Highlighting}
\end{Shaded}

\emph{(Full implementation: \texttt{report1\_scripts/qaoa\_circuit.py})}

\subsection{\texorpdfstring{Module 3:
\texttt{optimizer.py}}{Module 3: optimizer.py}}\label{module-3-optimizer.py}

\begin{Shaded}
\begin{Highlighting}[]
\CommentTok{\# {-}{-}{-} optimizer.py pseudocode {-}{-}{-}}

\KeywordTok{def}\NormalTok{ optimize\_qaoa(qnode, p, method, init\_params, maxiter, tol, callback):}
\NormalTok{    d }\OperatorTok{=} \DecValTok{2} \OperatorTok{*}\NormalTok{ p                         }\CommentTok{\# total number of parameters}
\NormalTok{    eval\_count }\OperatorTok{=} \DecValTok{0}                    \CommentTok{\# mutable circuit evaluation counter}
\NormalTok{    trajectory }\OperatorTok{=}\NormalTok{ []                   }\CommentTok{\# list of (params\_copy, energy)}

    \CommentTok{\# {-}{-}{-} Initialisation {-}{-}{-}}
    \ControlFlowTok{if}\NormalTok{ init\_params }\KeywordTok{is} \VariableTok{None}\NormalTok{:}
\NormalTok{        gammas }\OperatorTok{=}\NormalTok{ np.random.uniform(}\DecValTok{0}\NormalTok{, np.pi, size}\OperatorTok{=}\NormalTok{p)       }\CommentTok{\# shape (p,)}
\NormalTok{        betas }\OperatorTok{=}\NormalTok{ np.random.uniform(}\DecValTok{0}\NormalTok{, np.pi }\OperatorTok{/} \DecValTok{2}\NormalTok{, size}\OperatorTok{=}\NormalTok{p)    }\CommentTok{\# shape (p,)}
\NormalTok{        params }\OperatorTok{=}\NormalTok{ np.concatenate([gammas, betas])            }\CommentTok{\# shape (d,)}
    \ControlFlowTok{else}\NormalTok{:}
\NormalTok{        params }\OperatorTok{=}\NormalTok{ init\_params.copy()}

    \CommentTok{\# Negated objective: optimisers minimise, QAOA maximises \textless{}H\_C\textgreater{}}
    \KeywordTok{def}\NormalTok{ objective(theta):}
        \KeywordTok{nonlocal}\NormalTok{ eval\_count}
\NormalTok{        eval\_count }\OperatorTok{+=} \DecValTok{1}
\NormalTok{        val }\OperatorTok{=} \BuiltInTok{float}\NormalTok{(qnode(theta))}
        \ControlFlowTok{return} \OperatorTok{{-}}\NormalTok{val}

    \CommentTok{\# {-}{-}{-} COBYLA {-}{-}{-}}
    \ControlFlowTok{if}\NormalTok{ method }\OperatorTok{==} \StringTok{"COBYLA"}\NormalTok{:}
        \ImportTok{from}\NormalTok{ scipy.optimize }\ImportTok{import}\NormalTok{ minimize}
\NormalTok{        result }\OperatorTok{=}\NormalTok{ minimize(}
\NormalTok{            objective, params, method}\OperatorTok{=}\StringTok{"COBYLA"}\NormalTok{,}
\NormalTok{            options}\OperatorTok{=}\NormalTok{\{}\StringTok{"maxiter"}\NormalTok{: maxiter, }\StringTok{"rhobeg"}\NormalTok{: }\FloatTok{0.5}\NormalTok{\},}
\NormalTok{            callback}\OperatorTok{=}\KeywordTok{lambda}\NormalTok{ x: trajectory.append((x.copy(), }\OperatorTok{{-}}\NormalTok{objective(x))),}
\NormalTok{        )}
\NormalTok{        best\_params }\OperatorTok{=}\NormalTok{ result.x}

    \CommentTok{\# {-}{-}{-} L{-}BFGS{-}B {-}{-}{-}}
    \ControlFlowTok{elif}\NormalTok{ method }\OperatorTok{==} \StringTok{"L{-}BFGS{-}B"}\NormalTok{:}
        \ImportTok{from}\NormalTok{ scipy.optimize }\ImportTok{import}\NormalTok{ minimize}
        \KeywordTok{def}\NormalTok{ jac(theta):}
            \ControlFlowTok{return} \OperatorTok{{-}}\NormalTok{parameter\_shift\_gradient(qnode, theta)}
\NormalTok{        result }\OperatorTok{=}\NormalTok{ minimize(}
\NormalTok{            objective, params, method}\OperatorTok{=}\StringTok{"L{-}BFGS{-}B"}\NormalTok{, jac}\OperatorTok{=}\NormalTok{jac,}
\NormalTok{            options}\OperatorTok{=}\NormalTok{\{}\StringTok{"maxiter"}\NormalTok{: maxiter, }\StringTok{"ftol"}\NormalTok{: tol\},}
\NormalTok{        )}
\NormalTok{        best\_params }\OperatorTok{=}\NormalTok{ result.x}

    \CommentTok{\# {-}{-}{-} SPSA {-}{-}{-}}
    \ControlFlowTok{elif}\NormalTok{ method }\OperatorTok{==} \StringTok{"SPSA"}\NormalTok{:}
        \CommentTok{\# Gain schedule constants}
\NormalTok{        a0, c0 }\OperatorTok{=} \FloatTok{0.1}\NormalTok{, }\FloatTok{0.1}
\NormalTok{        alpha, gamma\_spsa }\OperatorTok{=} \FloatTok{0.602}\NormalTok{, }\FloatTok{0.101}
\NormalTok{        A }\OperatorTok{=} \FloatTok{0.1} \OperatorTok{*}\NormalTok{ maxiter}

        \ControlFlowTok{for}\NormalTok{ k }\KeywordTok{in} \BuiltInTok{range}\NormalTok{(}\DecValTok{1}\NormalTok{, maxiter }\OperatorTok{+} \DecValTok{1}\NormalTok{):}
\NormalTok{            a\_k }\OperatorTok{=}\NormalTok{ a0 }\OperatorTok{/}\NormalTok{ (k }\OperatorTok{+}\NormalTok{ A) }\OperatorTok{**}\NormalTok{ alpha       }\CommentTok{\# step size schedule}
\NormalTok{            c\_k }\OperatorTok{=}\NormalTok{ c0 }\OperatorTok{/}\NormalTok{ k }\OperatorTok{**}\NormalTok{ gamma\_spsa         }\CommentTok{\# perturbation size schedule}

            \CommentTok{\# Rademacher random perturbation vector: each entry +1 or {-}1}
\NormalTok{            delta }\OperatorTok{=}\NormalTok{ np.random.choice([}\OperatorTok{{-}}\DecValTok{1}\NormalTok{, }\DecValTok{1}\NormalTok{], size}\OperatorTok{=}\NormalTok{d)   }\CommentTok{\# shape (d,)}

            \CommentTok{\# Two{-}point gradient estimate (only 2 evaluations per step)}
\NormalTok{            f\_plus }\OperatorTok{=}\NormalTok{ objective(params }\OperatorTok{+}\NormalTok{ c\_k }\OperatorTok{*}\NormalTok{ delta)}
\NormalTok{            f\_minus }\OperatorTok{=}\NormalTok{ objective(params }\OperatorTok{{-}}\NormalTok{ c\_k }\OperatorTok{*}\NormalTok{ delta)}
\NormalTok{            g\_hat }\OperatorTok{=}\NormalTok{ (f\_plus }\OperatorTok{{-}}\NormalTok{ f\_minus) }\OperatorTok{/}\NormalTok{ (}\DecValTok{2} \OperatorTok{*}\NormalTok{ c\_k) }\OperatorTok{*}\NormalTok{ (}\FloatTok{1.0} \OperatorTok{/}\NormalTok{ delta)  }\CommentTok{\# shape (d,)}

\NormalTok{            params }\OperatorTok{=}\NormalTok{ params }\OperatorTok{{-}}\NormalTok{ a\_k }\OperatorTok{*}\NormalTok{ g\_hat      }\CommentTok{\# gradient descent step}
\NormalTok{            trajectory.append((params.copy(), }\OperatorTok{{-}}\BuiltInTok{float}\NormalTok{(qnode(params))))}

            \ControlFlowTok{if}\NormalTok{ callback:}
\NormalTok{                callback(params, trajectory[}\OperatorTok{{-}}\DecValTok{1}\NormalTok{][}\DecValTok{1}\NormalTok{], k)}

\NormalTok{        best\_params }\OperatorTok{=}\NormalTok{ params}

    \CommentTok{\# {-}{-}{-} Adam {-}{-}{-}}
    \ControlFlowTok{elif}\NormalTok{ method }\OperatorTok{==} \StringTok{"Adam"}\NormalTok{:}
\NormalTok{        opt }\OperatorTok{=}\NormalTok{ qml.AdamOptimizer(stepsize}\OperatorTok{=}\FloatTok{0.05}\NormalTok{)}
        \ControlFlowTok{for}\NormalTok{ k }\KeywordTok{in} \BuiltInTok{range}\NormalTok{(maxiter):}
\NormalTok{            params, cost }\OperatorTok{=}\NormalTok{ opt.step\_and\_cost(}\KeywordTok{lambda}\NormalTok{ t: }\OperatorTok{{-}}\NormalTok{qnode(t), params)}
\NormalTok{            energy }\OperatorTok{=} \OperatorTok{{-}}\NormalTok{cost}
\NormalTok{            trajectory.append((params.copy(), energy))}
            \ControlFlowTok{if}\NormalTok{ callback:}
\NormalTok{                callback(params, energy, k)}
            \CommentTok{\# Convergence check}
            \ControlFlowTok{if}\NormalTok{ k }\OperatorTok{\textgreater{}} \DecValTok{0} \KeywordTok{and} \BuiltInTok{abs}\NormalTok{(trajectory[}\OperatorTok{{-}}\DecValTok{1}\NormalTok{][}\DecValTok{1}\NormalTok{] }\OperatorTok{{-}}\NormalTok{ trajectory[}\OperatorTok{{-}}\DecValTok{2}\NormalTok{][}\DecValTok{1}\NormalTok{]) }\OperatorTok{\textless{}}\NormalTok{ tol:}
                \ControlFlowTok{break}
\NormalTok{        best\_params }\OperatorTok{=}\NormalTok{ params}

    \ControlFlowTok{return}\NormalTok{ OptimizationResult(}
\NormalTok{        best\_params}\OperatorTok{=}\NormalTok{best\_params,}
\NormalTok{        optimal\_energy}\OperatorTok{={-}}\NormalTok{objective(best\_params),}
\NormalTok{        trajectory}\OperatorTok{=}\NormalTok{trajectory,}
\NormalTok{        n\_evaluations}\OperatorTok{=}\NormalTok{eval\_count,}
\NormalTok{        converged}\OperatorTok{=}\NormalTok{...,}
\NormalTok{    )}


\KeywordTok{def}\NormalTok{ parameter\_shift\_gradient(qnode, params, shift}\OperatorTok{=}\NormalTok{np.pi }\OperatorTok{/} \DecValTok{2}\NormalTok{):}
\NormalTok{    d }\OperatorTok{=} \BuiltInTok{len}\NormalTok{(params)                    }\CommentTok{\# number of parameters}
\NormalTok{    grad }\OperatorTok{=}\NormalTok{ np.zeros(d)                 }\CommentTok{\# shape (d,)}
    \ControlFlowTok{for}\NormalTok{ i }\KeywordTok{in} \BuiltInTok{range}\NormalTok{(d):}
\NormalTok{        shifted\_plus }\OperatorTok{=}\NormalTok{ params.copy()}
\NormalTok{        shifted\_minus }\OperatorTok{=}\NormalTok{ params.copy()}
\NormalTok{        shifted\_plus[i] }\OperatorTok{+=}\NormalTok{ shift       }\CommentTok{\# theta\_i + pi/2}
\NormalTok{        shifted\_minus[i] }\OperatorTok{{-}=}\NormalTok{ shift      }\CommentTok{\# theta\_i {-} pi/2}
        \CommentTok{\# Exact parameter{-}shift formula (valid for eigenvalues +/{-}1)}
\NormalTok{        grad[i] }\OperatorTok{=}\NormalTok{ (}\BuiltInTok{float}\NormalTok{(qnode(shifted\_plus)) }\OperatorTok{{-}} \BuiltInTok{float}\NormalTok{(qnode(shifted\_minus))) }\OperatorTok{/} \FloatTok{2.0}
    \CommentTok{\# Total: 2*d circuit evaluations}
    \ControlFlowTok{return}\NormalTok{ grad}
\end{Highlighting}
\end{Shaded}

\emph{(Full implementation: \texttt{report1\_scripts/optimizer.py})}

\subsection{\texorpdfstring{Module 4:
\texttt{landscape\_sampler.py}}{Module 4: landscape\_sampler.py}}\label{module-4-landscape_sampler.py}

\begin{Shaded}
\begin{Highlighting}[]
\CommentTok{\# {-}{-}{-} landscape\_sampler.py pseudocode {-}{-}{-}}

\KeywordTok{def}\NormalTok{ sample\_landscape(qnode, p, gamma\_range, beta\_range, n\_gamma, n\_beta,}
\NormalTok{                     fixed\_params, c\_max):}
\NormalTok{    gamma\_vals }\OperatorTok{=}\NormalTok{ np.linspace(gamma\_range[}\DecValTok{0}\NormalTok{], gamma\_range[}\DecValTok{1}\NormalTok{], n\_gamma)  }\CommentTok{\# (N\_gamma,)}
\NormalTok{    beta\_vals }\OperatorTok{=}\NormalTok{ np.linspace(beta\_range[}\DecValTok{0}\NormalTok{], beta\_range[}\DecValTok{1}\NormalTok{], n\_beta)      }\CommentTok{\# (N\_beta,)}
\NormalTok{    energies }\OperatorTok{=}\NormalTok{ np.zeros((n\_gamma, n\_beta))                              }\CommentTok{\# (N\_gamma, N\_beta)}

    \ControlFlowTok{for}\NormalTok{ i, gamma }\KeywordTok{in} \BuiltInTok{enumerate}\NormalTok{(gamma\_vals):}
        \ControlFlowTok{for}\NormalTok{ j, beta }\KeywordTok{in} \BuiltInTok{enumerate}\NormalTok{(beta\_vals):}
            \ControlFlowTok{if}\NormalTok{ p }\OperatorTok{==} \DecValTok{1}\NormalTok{:}
\NormalTok{                params }\OperatorTok{=}\NormalTok{ np.array([gamma, beta])                        }\CommentTok{\# (2,)}
            \ControlFlowTok{else}\NormalTok{:}
                \CommentTok{\# 2D slice through higher{-}dimensional landscape}
\NormalTok{                params }\OperatorTok{=}\NormalTok{ fixed\_params.copy()                            }\CommentTok{\# (2p,)}
\NormalTok{                params[}\DecValTok{0}\NormalTok{] }\OperatorTok{=}\NormalTok{ gamma      }\CommentTok{\# overwrite first gamma}
\NormalTok{                params[p] }\OperatorTok{=}\NormalTok{ beta       }\CommentTok{\# overwrite first beta}
\NormalTok{            energies[i, j] }\OperatorTok{=} \BuiltInTok{float}\NormalTok{(qnode(params))}

\NormalTok{    approx\_ratios }\OperatorTok{=}\NormalTok{ energies }\OperatorTok{/}\NormalTok{ c\_max                                    }\CommentTok{\# (N\_gamma, N\_beta)}
    \ControlFlowTok{return}\NormalTok{ LandscapeData(gamma\_vals, beta\_vals, energies, approx\_ratios)}


\KeywordTok{def}\NormalTok{ sample\_landscape\_parallel(qnode, p, gamma\_range, beta\_range,}
\NormalTok{                              n\_gamma, n\_beta, c\_max):}
\NormalTok{    gamma\_vals }\OperatorTok{=}\NormalTok{ np.linspace(gamma\_range[}\DecValTok{0}\NormalTok{], gamma\_range[}\DecValTok{1}\NormalTok{], n\_gamma)  }\CommentTok{\# (N\_gamma,)}
\NormalTok{    beta\_vals }\OperatorTok{=}\NormalTok{ np.linspace(beta\_range[}\DecValTok{0}\NormalTok{], beta\_range[}\DecValTok{1}\NormalTok{], n\_beta)      }\CommentTok{\# (N\_beta,)}

    \CommentTok{\# Build full parameter grid via meshgrid with \textquotesingle{}ij\textquotesingle{} indexing}
\NormalTok{    gamma\_grid, beta\_grid }\OperatorTok{=}\NormalTok{ np.meshgrid(gamma\_vals, beta\_vals, indexing}\OperatorTok{=}\StringTok{\textquotesingle{}ij\textquotesingle{}}\NormalTok{)}
    \CommentTok{\# gamma\_grid shape: (N\_gamma, N\_beta)}
    \CommentTok{\# beta\_grid  shape: (N\_gamma, N\_beta)}

    \CommentTok{\# Flatten into 1D arrays of length N\_gamma * N\_beta}
\NormalTok{    gamma\_flat }\OperatorTok{=}\NormalTok{ gamma\_grid.ravel()                                     }\CommentTok{\# (N\_gamma*N\_beta,)}
\NormalTok{    beta\_flat }\OperatorTok{=}\NormalTok{ beta\_grid.ravel()                                       }\CommentTok{\# (N\_gamma*N\_beta,)}

    \CommentTok{\# Assemble parameter matrix for all grid points}
\NormalTok{    n\_points }\OperatorTok{=}\NormalTok{ n\_gamma }\OperatorTok{*}\NormalTok{ n\_beta}
\NormalTok{    all\_params }\OperatorTok{=}\NormalTok{ np.zeros((n\_points, }\DecValTok{2} \OperatorTok{*}\NormalTok{ p))                            }\CommentTok{\# (N\_total, 2p)}
\NormalTok{    all\_params[:, }\DecValTok{0}\NormalTok{] }\OperatorTok{=}\NormalTok{ gamma\_flat      }\CommentTok{\# first gamma for each point}
\NormalTok{    all\_params[:, p] }\OperatorTok{=}\NormalTok{ beta\_flat       }\CommentTok{\# first beta for each point}

    \CommentTok{\# Convert each parameter set into a PennyLane QuantumTape}
\NormalTok{    tapes }\OperatorTok{=}\NormalTok{ []}
    \ControlFlowTok{for}\NormalTok{ idx }\KeywordTok{in} \BuiltInTok{range}\NormalTok{(n\_points):}
\NormalTok{        tape }\OperatorTok{=}\NormalTok{ qml.tape.QuantumScript(...)  }\CommentTok{\# build tape from qnode with all\_params[idx]}
\NormalTok{        tapes.append(tape)}

    \CommentTok{\# Execute all tapes in one batch {-}{-} exploits simulator vectorisation}
\NormalTok{    results }\OperatorTok{=}\NormalTok{ qml.execute(tapes, device}\OperatorTok{=}\NormalTok{qnode.device, gradient\_fn}\OperatorTok{=}\VariableTok{None}\NormalTok{)}
    \CommentTok{\# results: list of length N\_total, each a scalar expectation value}

\NormalTok{    energies }\OperatorTok{=}\NormalTok{ np.array(results).reshape(n\_gamma, n\_beta)               }\CommentTok{\# (N\_gamma, N\_beta)}
\NormalTok{    approx\_ratios }\OperatorTok{=}\NormalTok{ energies }\OperatorTok{/}\NormalTok{ c\_max}
    \ControlFlowTok{return}\NormalTok{ LandscapeData(gamma\_vals, beta\_vals, energies, approx\_ratios)}
\end{Highlighting}
\end{Shaded}

\emph{(Full implementation:
\texttt{report1\_scripts/landscape\_sampler.py})}

\subsection{\texorpdfstring{Module 5:
\texttt{noise\_model.py}}{Module 5: noise\_model.py}}\label{module-5-noise_model.py}

\begin{Shaded}
\begin{Highlighting}[]
\CommentTok{\# {-}{-}{-} noise\_model.py pseudocode {-}{-}{-}}

\KeywordTok{def}\NormalTok{ create\_noisy\_device(n\_qubits, }\OperatorTok{**}\NormalTok{noise\_kwargs):}
    \CommentTok{\# Device is noise{-}agnostic; channels are inserted inside the QNode}
    \ControlFlowTok{return}\NormalTok{ qml.device(}\StringTok{"default.mixed"}\NormalTok{, wires}\OperatorTok{=}\NormalTok{n\_qubits)}


\KeywordTok{def}\NormalTok{ noisy\_qaoa\_qnode(cost\_hamiltonian, n\_qubits, p, noise\_params):}
\NormalTok{    dev }\OperatorTok{=}\NormalTok{ create\_noisy\_device(n\_qubits, }\OperatorTok{**}\NormalTok{noise\_params)}
\NormalTok{    edges }\OperatorTok{=}\NormalTok{ extract\_edges\_from\_hamiltonian(cost\_hamiltonian)}

    \AttributeTok{@qml.qnode}\NormalTok{(dev, diff\_method}\OperatorTok{=}\StringTok{"best"}\NormalTok{)}
    \KeywordTok{def}\NormalTok{ circuit(params):}
\NormalTok{        gammas }\OperatorTok{=}\NormalTok{ params[:p]                                   }\CommentTok{\# shape (p,)}
\NormalTok{        betas }\OperatorTok{=}\NormalTok{ params[p:]                                    }\CommentTok{\# shape (p,)}

        \CommentTok{\# {-}{-}{-} Hadamard layer + noise {-}{-}{-}}
        \ControlFlowTok{for}\NormalTok{ q }\KeywordTok{in} \BuiltInTok{range}\NormalTok{(n\_qubits):}
\NormalTok{            qml.Hadamard(wires}\OperatorTok{=}\NormalTok{q)}
            \CommentTok{\# Insert single{-}qubit noise after each Hadamard}
\NormalTok{            qml.DepolarizingChannel(noise\_params[}\StringTok{"depol\_rate"}\NormalTok{], wires}\OperatorTok{=}\NormalTok{q)}
\NormalTok{            qml.ThermalRelaxationError(}
\NormalTok{                noise\_params[}\StringTok{"t1"}\NormalTok{], noise\_params[}\StringTok{"t2"}\NormalTok{],}
\NormalTok{                noise\_params[}\StringTok{"gate\_time\_1q"}\NormalTok{], wires}\OperatorTok{=}\NormalTok{q}
\NormalTok{            )}

        \ControlFlowTok{for}\NormalTok{ k }\KeywordTok{in} \BuiltInTok{range}\NormalTok{(p):}
            \CommentTok{\# {-}{-}{-} Cost layer + two{-}qubit noise {-}{-}{-}}
            \ControlFlowTok{for}\NormalTok{ i, j, w }\KeywordTok{in}\NormalTok{ edges:}
\NormalTok{                qml.IsingZZ(gammas[k] }\OperatorTok{*}\NormalTok{ w, wires}\OperatorTok{=}\NormalTok{[i, j])}
                \CommentTok{\# Two{-}qubit noise: \textasciitilde{}10x single{-}qubit depol rate}
                \ControlFlowTok{for}\NormalTok{ qubit }\KeywordTok{in}\NormalTok{ [i, j]:}
\NormalTok{                    qml.DepolarizingChannel(}
\NormalTok{                        noise\_params[}\StringTok{"depol\_rate"}\NormalTok{] }\OperatorTok{*} \DecValTok{10}\NormalTok{, wires}\OperatorTok{=}\NormalTok{qubit}
\NormalTok{                    )}
\NormalTok{                    qml.ThermalRelaxationError(}
\NormalTok{                        noise\_params[}\StringTok{"t1"}\NormalTok{], noise\_params[}\StringTok{"t2"}\NormalTok{],}
\NormalTok{                        noise\_params[}\StringTok{"gate\_time\_2q"}\NormalTok{], wires}\OperatorTok{=}\NormalTok{qubit}
\NormalTok{                    )}

            \CommentTok{\# {-}{-}{-} Mixer layer + single{-}qubit noise {-}{-}{-}}
            \ControlFlowTok{for}\NormalTok{ q }\KeywordTok{in} \BuiltInTok{range}\NormalTok{(n\_qubits):}
\NormalTok{                qml.RX(}\DecValTok{2} \OperatorTok{*}\NormalTok{ betas[k], wires}\OperatorTok{=}\NormalTok{q)}
\NormalTok{                qml.DepolarizingChannel(noise\_params[}\StringTok{"depol\_rate"}\NormalTok{], wires}\OperatorTok{=}\NormalTok{q)}
\NormalTok{                qml.ThermalRelaxationError(}
\NormalTok{                    noise\_params[}\StringTok{"t1"}\NormalTok{], noise\_params[}\StringTok{"t2"}\NormalTok{],}
\NormalTok{                    noise\_params[}\StringTok{"gate\_time\_1q"}\NormalTok{], wires}\OperatorTok{=}\NormalTok{q}
\NormalTok{                )}

        \ControlFlowTok{return}\NormalTok{ qml.expval(cost\_hamiltonian)}
    \ControlFlowTok{return}\NormalTok{ circuit}


\KeywordTok{def}\NormalTok{ lindblad\_simulation(initial\_state, hamiltonian, collapse\_ops, tlist,}
\NormalTok{                        observables):}
    \CommentTok{\# initial\_state: qutip.Qobj density matrix of shape (2\^{}N, 2\^{}N)}
    \CommentTok{\# hamiltonian:   qutip.Qobj Hamiltonian of shape (2\^{}N, 2\^{}N)}
    \CommentTok{\# collapse\_ops:  list of qutip.Qobj, each shape (2\^{}N, 2\^{}N)}
    \CommentTok{\# tlist:         np.ndarray of time points, shape (T,)}

    \CommentTok{\# Build per{-}qubit collapse operators as tensor products:}
    \CommentTok{\#   Amplitude damping: L\_ad = sqrt(1/T1) * sigma\_minus, tensored}
    \CommentTok{\#   Pure dephasing:    L\_pd = sqrt(1/T2 {-} 1/(2*T1)) * Z/2, tensored}
    \ControlFlowTok{for}\NormalTok{ q }\KeywordTok{in} \BuiltInTok{range}\NormalTok{(n\_qubits):}
\NormalTok{        gamma\_1 }\OperatorTok{=} \FloatTok{1.0} \OperatorTok{/}\NormalTok{ T1}
\NormalTok{        gamma\_phi }\OperatorTok{=} \FloatTok{1.0} \OperatorTok{/}\NormalTok{ T2 }\OperatorTok{{-}} \FloatTok{1.0} \OperatorTok{/}\NormalTok{ (}\DecValTok{2} \OperatorTok{*}\NormalTok{ T1)}
        \CommentTok{\# Tensor product: I (x) ... (x) sigma\_minus (x) ... (x) I}
\NormalTok{        L\_ad }\OperatorTok{=}\NormalTok{ np.sqrt(gamma\_1) }\OperatorTok{*}\NormalTok{ tensor\_single\_op(qutip.destroy(}\DecValTok{2}\NormalTok{), q, n\_qubits)}
\NormalTok{        L\_pd }\OperatorTok{=}\NormalTok{ np.sqrt(gamma\_phi) }\OperatorTok{*}\NormalTok{ tensor\_single\_op(qutip.sigmaz() }\OperatorTok{/} \DecValTok{2}\NormalTok{, q, n\_qubits)}
\NormalTok{        collapse\_ops.extend([L\_ad, L\_pd])}

    \CommentTok{\# Solve Lindblad master equation}
\NormalTok{    result }\OperatorTok{=}\NormalTok{ qutip.mesolve(}
\NormalTok{        hamiltonian, initial\_state, tlist,}
\NormalTok{        c\_ops}\OperatorTok{=}\NormalTok{collapse\_ops,}
\NormalTok{        e\_ops}\OperatorTok{=}\NormalTok{[observables[label] }\ControlFlowTok{for}\NormalTok{ label }\KeywordTok{in}\NormalTok{ observables],}
\NormalTok{    )}
    \CommentTok{\# result.expect[k] is np.ndarray of shape (T,) for k{-}th observable}
    \ControlFlowTok{return}\NormalTok{ \{label: result.expect[k] }\ControlFlowTok{for}\NormalTok{ k, label }\KeywordTok{in} \BuiltInTok{enumerate}\NormalTok{(observables)\}}


\KeywordTok{def}\NormalTok{ compare\_noise\_channels(n\_qubits, circuit\_params, noise\_params):}
    \CommentTok{\# 1. PennyLane channel model}
\NormalTok{    qnode\_pl }\OperatorTok{=}\NormalTok{ noisy\_qaoa\_qnode(...)}
\NormalTok{    energy\_pl }\OperatorTok{=} \BuiltInTok{float}\NormalTok{(qnode\_pl(circuit\_params))}
    \CommentTok{\# Extract density matrix from PennyLane (via qml.state())}
\NormalTok{    rho\_pl }\OperatorTok{=}\NormalTok{ ...                         }\CommentTok{\# shape (2\^{}N, 2\^{}N) complex}

    \CommentTok{\# 2. QuTiP Lindblad model}
    \CommentTok{\# Initialise |+\textgreater{}\^{}N as a density matrix}
\NormalTok{    plus }\OperatorTok{=}\NormalTok{ (qutip.basis(}\DecValTok{2}\NormalTok{, }\DecValTok{0}\NormalTok{) }\OperatorTok{+}\NormalTok{ qutip.basis(}\DecValTok{2}\NormalTok{, }\DecValTok{1}\NormalTok{)).unit()}
\NormalTok{    rho0 }\OperatorTok{=}\NormalTok{ qutip.tensor([plus }\OperatorTok{*}\NormalTok{ plus.dag()] }\OperatorTok{*}\NormalTok{ n\_qubits)    }\CommentTok{\# (2\^{}N, 2\^{}N)}
    \CommentTok{\# Step through each gate as piecewise{-}constant Hamiltonian evolution}
    \CommentTok{\# ... (build H\_gate for each gate, evolve for gate\_time) ...}
\NormalTok{    result }\OperatorTok{=}\NormalTok{ lindblad\_simulation(rho0, H\_total, collapse\_ops, tlist, ...)}
\NormalTok{    energy\_qt }\OperatorTok{=}\NormalTok{ result[}\StringTok{"H\_C"}\NormalTok{][}\OperatorTok{{-}}\DecValTok{1}\NormalTok{]        }\CommentTok{\# final time{-}step energy}

    \CommentTok{\# 3. Compare}
\NormalTok{    fidelity }\OperatorTok{=}\NormalTok{ qutip.fidelity(rho\_pl\_qobj, rho\_qt)}
\NormalTok{    purity\_pl }\OperatorTok{=}\NormalTok{ np.real(np.trace(rho\_pl }\OperatorTok{@}\NormalTok{ rho\_pl))}
\NormalTok{    purity\_qt }\OperatorTok{=}\NormalTok{ np.real((rho\_qt }\OperatorTok{*}\NormalTok{ rho\_qt).tr())}
    \ControlFlowTok{return}\NormalTok{ \{}\StringTok{"energy\_pennylane"}\NormalTok{: energy\_pl, }\StringTok{"energy\_qutip"}\NormalTok{: energy\_qt,}
            \StringTok{"fidelity"}\NormalTok{: fidelity, }\StringTok{"purity\_pl"}\NormalTok{: purity\_pl, }\StringTok{"purity\_qt"}\NormalTok{: purity\_qt\}}
\end{Highlighting}
\end{Shaded}

\emph{(Full implementation: \texttt{report1\_scripts/noise\_model.py})}

\subsection{\texorpdfstring{Module 6:
\texttt{hardware\_runner.py}}{Module 6: hardware\_runner.py}}\label{module-6-hardware_runner.py}

\begin{Shaded}
\begin{Highlighting}[]
\CommentTok{\# {-}{-}{-} hardware\_runner.py pseudocode {-}{-}{-}}

\KeywordTok{def}\NormalTok{ create\_ibm\_device(n\_qubits, backend\_name, shots, ibm\_token):}
    \ImportTok{import}\NormalTok{ os}
\NormalTok{    token }\OperatorTok{=}\NormalTok{ ibm\_token }\KeywordTok{or}\NormalTok{ os.environ.get(}\StringTok{"IBMQ\_TOKEN"}\NormalTok{)}
    \CommentTok{\# PennyLane{-}Qiskit plugin handles transpilation internally}
\NormalTok{    dev }\OperatorTok{=}\NormalTok{ qml.device(}\StringTok{"qiskit.ibmq"}\NormalTok{, wires}\OperatorTok{=}\NormalTok{n\_qubits,}
\NormalTok{                      backend}\OperatorTok{=}\NormalTok{backend\_name, shots}\OperatorTok{=}\NormalTok{shots, ibmqx\_token}\OperatorTok{=}\NormalTok{token)}
    \ControlFlowTok{return}\NormalTok{ dev}


\KeywordTok{def}\NormalTok{ create\_iqm\_device(n\_qubits, backend\_url, shots):}
\NormalTok{    dev }\OperatorTok{=}\NormalTok{ qml.device(}\StringTok{"iqm.iqm\_device"}\NormalTok{, wires}\OperatorTok{=}\NormalTok{n\_qubits,}
\NormalTok{                      url}\OperatorTok{=}\NormalTok{backend\_url, shots}\OperatorTok{=}\NormalTok{shots)}
    \ControlFlowTok{return}\NormalTok{ dev}


\KeywordTok{def}\NormalTok{ execute\_on\_hardware(qnode, params, cost\_hamiltonian):}
    \ImportTok{import}\NormalTok{ time}
\NormalTok{    t0 }\OperatorTok{=}\NormalTok{ time.perf\_counter()}

    \CommentTok{\# Strategy 1: direct expectation (preferred)}
\NormalTok{    energy }\OperatorTok{=} \BuiltInTok{float}\NormalTok{(qnode(params))         }\CommentTok{\# qml.expval(H\_C) internally}

    \CommentTok{\# Strategy 2 (alternative): sample{-}based post{-}processing}
    \CommentTok{\# samples = qnode\_sample(params)      \# shape (S, N) binary array}
    \CommentTok{\# for s in range(S):}
    \CommentTok{\#     z = samples[s]                  \# shape (N,) binary}
    \CommentTok{\#     spin = 1 {-} 2 * z               \# \{0,1\} {-}\textgreater{} \{+1,{-}1\}}
    \CommentTok{\#     for i, j, w in edges:}
    \CommentTok{\#         C\_s += w/2 * (1 {-} spin[i] * spin[j])}
    \CommentTok{\# energy = np.mean(C\_values)}

\NormalTok{    wall\_time }\OperatorTok{=}\NormalTok{ time.perf\_counter() }\OperatorTok{{-}}\NormalTok{ t0}
    \ControlFlowTok{return}\NormalTok{ HardwareResult(energy, ..., shots, backend\_name, wall\_time)}


\KeywordTok{def}\NormalTok{ select\_optimal\_qubits(backend\_name, n\_qubits, ibm\_token):}
    \ImportTok{from}\NormalTok{ qiskit\_ibm\_runtime }\ImportTok{import}\NormalTok{ QiskitRuntimeService}
\NormalTok{    service }\OperatorTok{=}\NormalTok{ QiskitRuntimeService(channel}\OperatorTok{=}\StringTok{"ibm\_quantum"}\NormalTok{, token}\OperatorTok{=}\NormalTok{ibm\_token)}
\NormalTok{    backend }\OperatorTok{=}\NormalTok{ service.backend(backend\_name)}
\NormalTok{    props }\OperatorTok{=}\NormalTok{ backend.properties()}

    \CommentTok{\# Build weighted graph over physical qubits}
\NormalTok{    coupling }\OperatorTok{=}\NormalTok{ backend.configuration().coupling\_map   }\CommentTok{\# list of [q\_i, q\_j] pairs}
\NormalTok{    G }\OperatorTok{=}\NormalTok{ nx.Graph()}
    \ControlFlowTok{for}\NormalTok{ qi, qj }\KeywordTok{in}\NormalTok{ coupling:}
\NormalTok{        cx\_error }\OperatorTok{=}\NormalTok{ props.gate\_error(}\StringTok{"cx"}\NormalTok{, [qi, qj])}
\NormalTok{        G.add\_edge(qi, qj, weight}\OperatorTok{=}\FloatTok{1.0} \OperatorTok{/}\NormalTok{ cx\_error)    }\CommentTok{\# weight = 1 / eps\_CX}

    \ControlFlowTok{for}\NormalTok{ q }\KeywordTok{in}\NormalTok{ G.nodes():}
\NormalTok{        G.nodes[q][}\StringTok{"t1"}\NormalTok{] }\OperatorTok{=}\NormalTok{ props.t1(q)}
\NormalTok{        G.nodes[q][}\StringTok{"t2"}\NormalTok{] }\OperatorTok{=}\NormalTok{ props.t2(q)}

    \CommentTok{\# Search for the best size{-}N connected subgraph via BFS enumeration}
\NormalTok{    best\_score }\OperatorTok{=} \OperatorTok{{-}}\NormalTok{np.inf}
\NormalTok{    best\_subgraph }\OperatorTok{=} \VariableTok{None}
\NormalTok{    alpha }\OperatorTok{=} \FloatTok{0.1}
    \ControlFlowTok{for}\NormalTok{ seed\_node }\KeywordTok{in}\NormalTok{ G.nodes():}
\NormalTok{        subgraph }\OperatorTok{=}\NormalTok{ bfs\_subgraph(G, seed\_node, n\_qubits)   }\CommentTok{\# first N reachable nodes}
\NormalTok{        score }\OperatorTok{=}\NormalTok{ (}
            \BuiltInTok{sum}\NormalTok{(G[u][v][}\StringTok{"weight"}\NormalTok{] }\ControlFlowTok{for}\NormalTok{ u, v }\KeywordTok{in}\NormalTok{ subgraph.edges())}
            \OperatorTok{+}\NormalTok{ alpha }\OperatorTok{*} \BuiltInTok{sum}\NormalTok{(G.nodes[q][}\StringTok{"t1"}\NormalTok{] }\OperatorTok{*}\NormalTok{ G.nodes[q][}\StringTok{"t2"}\NormalTok{]}
                          \ControlFlowTok{for}\NormalTok{ q }\KeywordTok{in}\NormalTok{ subgraph.nodes())}
\NormalTok{        )}
        \ControlFlowTok{if}\NormalTok{ score }\OperatorTok{\textgreater{}}\NormalTok{ best\_score:}
\NormalTok{            best\_score }\OperatorTok{=}\NormalTok{ score}
\NormalTok{            best\_subgraph }\OperatorTok{=}\NormalTok{ subgraph}
    \ControlFlowTok{return} \BuiltInTok{list}\NormalTok{(best\_subgraph.nodes())}
\end{Highlighting}
\end{Shaded}

\emph{(Full implementation:
\texttt{report1\_scripts/hardware\_runner.py})}

\subsection{\texorpdfstring{Module 7:
\texttt{error\_mitigation.py}}{Module 7: error\_mitigation.py}}\label{module-7-error_mitigation.py}

\begin{Shaded}
\begin{Highlighting}[]
\CommentTok{\# {-}{-}{-} error\_mitigation.py pseudocode {-}{-}{-}}

\KeywordTok{def}\NormalTok{ global\_fold(circuit\_fn, scale\_factor):}
    \CommentTok{"""Produce a noise{-}amplified circuit by global unitary folding."""}
\NormalTok{    n }\OperatorTok{=} \BuiltInTok{int}\NormalTok{((scale\_factor }\OperatorTok{{-}} \DecValTok{1}\NormalTok{) }\OperatorTok{/} \DecValTok{2}\NormalTok{)           }\CommentTok{\# full folding repetitions}
\NormalTok{    frac }\OperatorTok{=}\NormalTok{ (scale\_factor }\OperatorTok{{-}} \DecValTok{1}\NormalTok{) }\OperatorTok{/} \DecValTok{2} \OperatorTok{{-}}\NormalTok{ n         }\CommentTok{\# fractional remainder}

    \CommentTok{\# Record original circuit operations from PennyLane tape}
\NormalTok{    original\_tape }\OperatorTok{=}\NormalTok{ qml.tape.QuantumScript(...)}
\NormalTok{    ops }\OperatorTok{=}\NormalTok{ original\_tape.operations              }\CommentTok{\# list of gate operations}
\NormalTok{    N\_gates }\OperatorTok{=} \BuiltInTok{len}\NormalTok{(ops)}

    \KeywordTok{def}\NormalTok{ folded\_circuit(params):}
        \CommentTok{\# 1. Apply original circuit U}
        \ControlFlowTok{for}\NormalTok{ op }\KeywordTok{in}\NormalTok{ ops:}
\NormalTok{            qml.}\BuiltInTok{apply}\NormalTok{(op)}

        \CommentTok{\# 2. Apply n full folds: (U\^{}dag U) repeated n times}
        \ControlFlowTok{for}\NormalTok{ \_ }\KeywordTok{in} \BuiltInTok{range}\NormalTok{(n):}
            \CommentTok{\# U\^{}dag: apply ops in reverse, each adjointed}
            \ControlFlowTok{for}\NormalTok{ op }\KeywordTok{in} \BuiltInTok{reversed}\NormalTok{(ops):}
\NormalTok{                qml.adjoint(op)}
            \CommentTok{\# U: apply ops forward again}
            \ControlFlowTok{for}\NormalTok{ op }\KeywordTok{in}\NormalTok{ ops:}
\NormalTok{                qml.}\BuiltInTok{apply}\NormalTok{(op)}

        \CommentTok{\# 3. Partial fold for non{-}integer scale factors}
        \ControlFlowTok{if}\NormalTok{ frac }\OperatorTok{\textgreater{}} \DecValTok{0}\NormalTok{:}
\NormalTok{            n\_partial }\OperatorTok{=} \BuiltInTok{int}\NormalTok{(}\BuiltInTok{round}\NormalTok{(frac }\OperatorTok{*}\NormalTok{ N\_gates))  }\CommentTok{\# gates to partially fold}
\NormalTok{            tail\_ops }\OperatorTok{=}\NormalTok{ ops[}\OperatorTok{{-}}\NormalTok{n\_partial:]              }\CommentTok{\# last n\_partial gates}
            \CommentTok{\# Adjoint of the tail}
            \ControlFlowTok{for}\NormalTok{ op }\KeywordTok{in} \BuiltInTok{reversed}\NormalTok{(tail\_ops):}
\NormalTok{                qml.adjoint(op)}
            \CommentTok{\# Re{-}apply the tail}
            \ControlFlowTok{for}\NormalTok{ op }\KeywordTok{in}\NormalTok{ tail\_ops:}
\NormalTok{                qml.}\BuiltInTok{apply}\NormalTok{(op)}

    \ControlFlowTok{return}\NormalTok{ folded\_circuit}


\KeywordTok{def}\NormalTok{ zne\_extrapolate(qnode, params, scale\_factors, extrapolation):}
\NormalTok{    noisy\_values }\OperatorTok{=}\NormalTok{ []                              }\CommentTok{\# shape will be (M,)}
    \ControlFlowTok{for}\NormalTok{ lam }\KeywordTok{in}\NormalTok{ scale\_factors:}
        \CommentTok{\# Build a folded QNode at scale factor lambda}
\NormalTok{        folded\_qnode }\OperatorTok{=}\NormalTok{ build\_folded\_qnode(qnode, lam)   }\CommentTok{\# uses global\_fold}
\NormalTok{        E\_lam }\OperatorTok{=} \BuiltInTok{float}\NormalTok{(folded\_qnode(params))}
\NormalTok{        noisy\_values.append(E\_lam)}

\NormalTok{    scale\_arr }\OperatorTok{=}\NormalTok{ np.array(scale\_factors)            }\CommentTok{\# shape (M,)}
\NormalTok{    noisy\_arr }\OperatorTok{=}\NormalTok{ np.array(noisy\_values)             }\CommentTok{\# shape (M,)}

    \ControlFlowTok{if}\NormalTok{ extrapolation }\OperatorTok{==} \StringTok{"linear"}\NormalTok{:}
\NormalTok{        E0 }\OperatorTok{=}\NormalTok{ linear\_extrapolation(scale\_arr, noisy\_arr)}
    \ControlFlowTok{elif}\NormalTok{ extrapolation }\OperatorTok{==} \StringTok{"richardson"}\NormalTok{:}
\NormalTok{        E0 }\OperatorTok{=}\NormalTok{ richardson\_extrapolation(scale\_arr, noisy\_arr)}

    \ControlFlowTok{return}\NormalTok{ E0, \{}\StringTok{"noisy\_values"}\NormalTok{: noisy\_arr, }\StringTok{"scale\_factors"}\NormalTok{: scale\_arr\}}


\KeywordTok{def}\NormalTok{ richardson\_extrapolation(scale\_factors, noisy\_values):}
    \CommentTok{\# scale\_factors: shape (M,)   noisy\_values: shape (M,)}
\NormalTok{    M }\OperatorTok{=} \BuiltInTok{len}\NormalTok{(scale\_factors)}
    \CommentTok{\# Build Vandermonde matrix V[i, j] = scale\_factors[i]\^{}j}
\NormalTok{    V }\OperatorTok{=}\NormalTok{ np.vander(scale\_factors, N}\OperatorTok{=}\NormalTok{M, increasing}\OperatorTok{=}\VariableTok{True}\NormalTok{)  }\CommentTok{\# shape (M, M)}
    \CommentTok{\# Solve V @ a = E for coefficient vector a}
\NormalTok{    a }\OperatorTok{=}\NormalTok{ np.linalg.solve(V, noisy\_values)                }\CommentTok{\# shape (M,)}
    \ControlFlowTok{return}\NormalTok{ a[}\DecValTok{0}\NormalTok{]                           }\CommentTok{\# a\_0 = E(lambda=0)}


\KeywordTok{def}\NormalTok{ linear\_extrapolation(scale\_factors, noisy\_values):}
    \CommentTok{\# Fit E(lambda) = a + b*lambda via least{-}squares}
    \CommentTok{\# Design matrix: [[1, lambda\_1], [1, lambda\_2], ...]}
\NormalTok{    A }\OperatorTok{=}\NormalTok{ np.column\_stack([np.ones\_like(scale\_factors), scale\_factors])  }\CommentTok{\# (M, 2)}
\NormalTok{    coeffs, }\OperatorTok{*}\NormalTok{\_ }\OperatorTok{=}\NormalTok{ np.linalg.lstsq(A, noisy\_values, rcond}\OperatorTok{=}\VariableTok{None}\NormalTok{)        }\CommentTok{\# (2,)}
    \ControlFlowTok{return}\NormalTok{ coeffs[}\DecValTok{0}\NormalTok{]                      }\CommentTok{\# a = E(0)}
\end{Highlighting}
\end{Shaded}

\emph{(Full implementation:
\texttt{report1\_scripts/error\_mitigation.py})}

\subsection{\texorpdfstring{Module 8:
\texttt{renderer.py}}{Module 8: renderer.py}}\label{module-8-renderer.py}

\begin{Shaded}
\begin{Highlighting}[]
\CommentTok{\# {-}{-}{-} renderer.py pseudocode {-}{-}{-}}

\KeywordTok{def}\NormalTok{ init\_opengl\_context(width, height, title):}
\NormalTok{    glfw.init()}
\NormalTok{    glfw.window\_hint(glfw.CONTEXT\_VERSION\_MAJOR, }\DecValTok{3}\NormalTok{)}
\NormalTok{    glfw.window\_hint(glfw.CONTEXT\_VERSION\_MINOR, }\DecValTok{3}\NormalTok{)}
\NormalTok{    glfw.window\_hint(glfw.OPENGL\_PROFILE, glfw.OPENGL\_CORE\_PROFILE)}
\NormalTok{    glfw.window\_hint(glfw.SAMPLES, }\DecValTok{4}\NormalTok{)              }\CommentTok{\# 4x MSAA}
\NormalTok{    window }\OperatorTok{=}\NormalTok{ glfw.create\_window(width, height, title, }\VariableTok{None}\NormalTok{, }\VariableTok{None}\NormalTok{)}
\NormalTok{    glfw.make\_context\_current(window)}
\NormalTok{    glEnable(GL\_DEPTH\_TEST)}
\NormalTok{    glEnable(GL\_MULTISAMPLE)}
\NormalTok{    glEnable(GL\_BLEND)}
\NormalTok{    glBlendFunc(GL\_SRC\_ALPHA, GL\_ONE\_MINUS\_SRC\_ALPHA)}
    \ControlFlowTok{return}\NormalTok{ window}


\KeywordTok{def}\NormalTok{ create\_surface\_vbo(gamma\_vals, beta\_vals, energies):}
\NormalTok{    n\_gamma, n\_beta }\OperatorTok{=} \BuiltInTok{len}\NormalTok{(gamma\_vals), }\BuiltInTok{len}\NormalTok{(beta\_vals)}

    \CommentTok{\# {-}{-}{-} Vertex data: 7 floats per vertex (28 bytes) {-}{-}{-}}
    \CommentTok{\# [gamma, beta, energy, nx, ny, nz, norm\_energy]}
\NormalTok{    vertices }\OperatorTok{=}\NormalTok{ np.zeros((n\_gamma }\OperatorTok{*}\NormalTok{ n\_beta, }\DecValTok{7}\NormalTok{), dtype}\OperatorTok{=}\NormalTok{np.float32)}
    \ControlFlowTok{for}\NormalTok{ i }\KeywordTok{in} \BuiltInTok{range}\NormalTok{(n\_gamma):}
        \ControlFlowTok{for}\NormalTok{ j }\KeywordTok{in} \BuiltInTok{range}\NormalTok{(n\_beta):}
\NormalTok{            idx }\OperatorTok{=}\NormalTok{ i }\OperatorTok{*}\NormalTok{ n\_beta }\OperatorTok{+}\NormalTok{ j}
\NormalTok{            E }\OperatorTok{=}\NormalTok{ energies[i, j]}
\NormalTok{            vertices[idx, }\DecValTok{0}\NormalTok{:}\DecValTok{3}\NormalTok{] }\OperatorTok{=}\NormalTok{ [gamma\_vals[i], beta\_vals[j], E]}
            \CommentTok{\# Finite{-}difference surface normals: n = ({-}dE/dgamma, {-}dE/dbeta, 1)}
\NormalTok{            dEdg }\OperatorTok{=}\NormalTok{ (energies[}\BuiltInTok{min}\NormalTok{(i}\OperatorTok{+}\DecValTok{1}\NormalTok{, n\_gamma}\OperatorTok{{-}}\DecValTok{1}\NormalTok{), j] }\OperatorTok{{-}}\NormalTok{ energies[}\BuiltInTok{max}\NormalTok{(i}\OperatorTok{{-}}\DecValTok{1}\NormalTok{, }\DecValTok{0}\NormalTok{), j]) }\OperatorTok{/}\NormalTok{ ...}
\NormalTok{            dEdb }\OperatorTok{=}\NormalTok{ (energies[i, }\BuiltInTok{min}\NormalTok{(j}\OperatorTok{+}\DecValTok{1}\NormalTok{, n\_beta}\OperatorTok{{-}}\DecValTok{1}\NormalTok{)] }\OperatorTok{{-}}\NormalTok{ energies[i, }\BuiltInTok{max}\NormalTok{(j}\OperatorTok{{-}}\DecValTok{1}\NormalTok{, }\DecValTok{0}\NormalTok{)]) }\OperatorTok{/}\NormalTok{ ...}
\NormalTok{            normal }\OperatorTok{=}\NormalTok{ np.array([}\OperatorTok{{-}}\NormalTok{dEdg, }\OperatorTok{{-}}\NormalTok{dEdb, }\FloatTok{1.0}\NormalTok{])}
\NormalTok{            normal }\OperatorTok{/=}\NormalTok{ np.linalg.norm(normal)}
\NormalTok{            vertices[idx, }\DecValTok{3}\NormalTok{:}\DecValTok{6}\NormalTok{] }\OperatorTok{=}\NormalTok{ normal}
\NormalTok{            vertices[idx, }\DecValTok{6}\NormalTok{] }\OperatorTok{=}\NormalTok{ ...                  }\CommentTok{\# normalised energy in [0, 1]}

    \CommentTok{\# {-}{-}{-} Index data: 2 triangles per grid cell {-}{-}{-}}
\NormalTok{    indices }\OperatorTok{=}\NormalTok{ []                                     }\CommentTok{\# uint32}
    \ControlFlowTok{for}\NormalTok{ i }\KeywordTok{in} \BuiltInTok{range}\NormalTok{(n\_gamma }\OperatorTok{{-}} \DecValTok{1}\NormalTok{):}
        \ControlFlowTok{for}\NormalTok{ j }\KeywordTok{in} \BuiltInTok{range}\NormalTok{(n\_beta }\OperatorTok{{-}} \DecValTok{1}\NormalTok{):}
\NormalTok{            tl }\OperatorTok{=}\NormalTok{ i }\OperatorTok{*}\NormalTok{ n\_beta }\OperatorTok{+}\NormalTok{ j                      }\CommentTok{\# top{-}left}
\NormalTok{            tr }\OperatorTok{=}\NormalTok{ (i }\OperatorTok{+} \DecValTok{1}\NormalTok{) }\OperatorTok{*}\NormalTok{ n\_beta }\OperatorTok{+}\NormalTok{ j                }\CommentTok{\# top{-}right}
\NormalTok{            bl }\OperatorTok{=}\NormalTok{ i }\OperatorTok{*}\NormalTok{ n\_beta }\OperatorTok{+}\NormalTok{ (j }\OperatorTok{+} \DecValTok{1}\NormalTok{)                }\CommentTok{\# bottom{-}left}
\NormalTok{            br }\OperatorTok{=}\NormalTok{ (i }\OperatorTok{+} \DecValTok{1}\NormalTok{) }\OperatorTok{*}\NormalTok{ n\_beta }\OperatorTok{+}\NormalTok{ (j }\OperatorTok{+} \DecValTok{1}\NormalTok{)          }\CommentTok{\# bottom{-}right}
\NormalTok{            indices.extend([tl, bl, tr, bl, br, tr])}
    \CommentTok{\# Total triangles: 2 * (N\_gamma {-} 1) * (N\_beta {-} 1)}

    \CommentTok{\# {-}{-}{-} GPU upload {-}{-}{-}}
\NormalTok{    vao }\OperatorTok{=}\NormalTok{ glGenVertexArrays(}\DecValTok{1}\NormalTok{)}
\NormalTok{    vbo }\OperatorTok{=}\NormalTok{ glGenBuffers(}\DecValTok{1}\NormalTok{)}
\NormalTok{    ebo }\OperatorTok{=}\NormalTok{ glGenBuffers(}\DecValTok{1}\NormalTok{)}
\NormalTok{    glBindVertexArray(vao)}
\NormalTok{    glBindBuffer(GL\_ARRAY\_BUFFER, vbo)}
\NormalTok{    glBufferData(GL\_ARRAY\_BUFFER, vertices.nbytes, vertices, GL\_DYNAMIC\_DRAW)}
\NormalTok{    glBindBuffer(GL\_ELEMENT\_ARRAY\_BUFFER, ebo)}
\NormalTok{    glBufferData(GL\_ELEMENT\_ARRAY\_BUFFER, ...)}
    \CommentTok{\# Attribute 0: position (3 floats, offset 0, stride 28)}
\NormalTok{    glVertexAttribPointer(}\DecValTok{0}\NormalTok{, }\DecValTok{3}\NormalTok{, GL\_FLOAT, GL\_FALSE, }\DecValTok{28}\NormalTok{, ctypes.c\_void\_p(}\DecValTok{0}\NormalTok{))}
    \CommentTok{\# Attribute 1: normal (3 floats, offset 12, stride 28)}
\NormalTok{    glVertexAttribPointer(}\DecValTok{1}\NormalTok{, }\DecValTok{3}\NormalTok{, GL\_FLOAT, GL\_FALSE, }\DecValTok{28}\NormalTok{, ctypes.c\_void\_p(}\DecValTok{12}\NormalTok{))}
    \CommentTok{\# Attribute 2: energy (1 float, offset 24, stride 28)}
\NormalTok{    glVertexAttribPointer(}\DecValTok{2}\NormalTok{, }\DecValTok{1}\NormalTok{, GL\_FLOAT, GL\_FALSE, }\DecValTok{28}\NormalTok{, ctypes.c\_void\_p(}\DecValTok{24}\NormalTok{))}
\NormalTok{    ...}
    \ControlFlowTok{return}\NormalTok{ vao, vbo, ebo, }\BuiltInTok{len}\NormalTok{(indices)}


\KeywordTok{def}\NormalTok{ draw\_optimizer\_trajectory(trajectory, shader\_program):}
\NormalTok{    n }\OperatorTok{=} \BuiltInTok{len}\NormalTok{(trajectory)}
    \CommentTok{\# Each vertex: (gamma, beta, energy, t) {-}{-} t in [0,1] for colour}
\NormalTok{    verts }\OperatorTok{=}\NormalTok{ np.zeros((n, }\DecValTok{4}\NormalTok{), dtype}\OperatorTok{=}\NormalTok{np.float32)}
    \ControlFlowTok{for}\NormalTok{ k, (gamma, beta, energy) }\KeywordTok{in} \BuiltInTok{enumerate}\NormalTok{(trajectory):}
\NormalTok{        verts[k] }\OperatorTok{=}\NormalTok{ [gamma, beta, energy }\OperatorTok{+} \FloatTok{0.01}\NormalTok{, k }\OperatorTok{/} \BuiltInTok{max}\NormalTok{(n }\OperatorTok{{-}} \DecValTok{1}\NormalTok{, }\DecValTok{1}\NormalTok{)]  }\CommentTok{\# z{-}offset}
    \CommentTok{\# Upload as GL\_STREAM\_DRAW, draw as GL\_LINE\_STRIP}
\NormalTok{    ...}
\end{Highlighting}
\end{Shaded}

\emph{(Full implementation: \texttt{report1\_scripts/renderer.py})}

\subsection{\texorpdfstring{Module 9:
\texttt{main.py}}{Module 9: main.py}}\label{module-9-main.py}

\begin{Shaded}
\begin{Highlighting}[]
\CommentTok{\# {-}{-}{-} main.py pseudocode {-}{-}{-}}

\KeywordTok{def}\NormalTok{ parse\_args():}
\NormalTok{    parser }\OperatorTok{=}\NormalTok{ argparse.ArgumentParser(description}\OperatorTok{=}\StringTok{"QAOA Energy Landscape Explorer"}\NormalTok{)}
\NormalTok{    parser.add\_argument(}\StringTok{"{-}{-}graph{-}type"}\NormalTok{, choices}\OperatorTok{=}\NormalTok{[...], default}\OperatorTok{=}\StringTok{"cycle"}\NormalTok{)}
\NormalTok{    parser.add\_argument(}\StringTok{"{-}{-}n{-}nodes"}\NormalTok{, }\BuiltInTok{type}\OperatorTok{=}\BuiltInTok{int}\NormalTok{, default}\OperatorTok{=}\DecValTok{6}\NormalTok{)}
\NormalTok{    parser.add\_argument(}\StringTok{"{-}{-}p"}\NormalTok{, }\BuiltInTok{type}\OperatorTok{=}\BuiltInTok{int}\NormalTok{, default}\OperatorTok{=}\DecValTok{1}\NormalTok{)}
\NormalTok{    parser.add\_argument(}\StringTok{"{-}{-}optimizer"}\NormalTok{, choices}\OperatorTok{=}\NormalTok{[}\StringTok{"COBYLA"}\NormalTok{,}\StringTok{"L{-}BFGS{-}B"}\NormalTok{,}\StringTok{"SPSA"}\NormalTok{,}\StringTok{"Adam"}\NormalTok{])}
\NormalTok{    parser.add\_argument(}\StringTok{"{-}{-}grid{-}size"}\NormalTok{, }\BuiltInTok{type}\OperatorTok{=}\BuiltInTok{int}\NormalTok{, default}\OperatorTok{=}\DecValTok{50}\NormalTok{)}
\NormalTok{    parser.add\_argument(}\StringTok{"{-}{-}noise"}\NormalTok{, action}\OperatorTok{=}\StringTok{"store\_true"}\NormalTok{)}
\NormalTok{    parser.add\_argument(}\StringTok{"{-}{-}hardware"}\NormalTok{, action}\OperatorTok{=}\StringTok{"store\_true"}\NormalTok{)}
\NormalTok{    parser.add\_argument(}\StringTok{"{-}{-}backend"}\NormalTok{, }\BuiltInTok{type}\OperatorTok{=}\BuiltInTok{str}\NormalTok{, default}\OperatorTok{=}\StringTok{"ibm\_brisbane"}\NormalTok{)}
\NormalTok{    parser.add\_argument(}\StringTok{"{-}{-}zne"}\NormalTok{, action}\OperatorTok{=}\StringTok{"store\_true"}\NormalTok{)}
\NormalTok{    parser.add\_argument(}\StringTok{"{-}{-}render"}\NormalTok{, action}\OperatorTok{=}\StringTok{"store\_true"}\NormalTok{)}
\NormalTok{    parser.add\_argument(}\StringTok{"{-}{-}config"}\NormalTok{, }\BuiltInTok{type}\OperatorTok{=}\BuiltInTok{str}\NormalTok{, default}\OperatorTok{=}\StringTok{"config.yaml"}\NormalTok{)}
\NormalTok{    parser.add\_argument(}\StringTok{"{-}{-}seed"}\NormalTok{, }\BuiltInTok{type}\OperatorTok{=}\BuiltInTok{int}\NormalTok{, default}\OperatorTok{=}\DecValTok{42}\NormalTok{)}
    \ControlFlowTok{return}\NormalTok{ parser.parse\_args()}


\KeywordTok{def}\NormalTok{ run\_pipeline(args):}
    \CommentTok{\# 1. Graph generation}
\NormalTok{    graph }\OperatorTok{=}\NormalTok{ create\_graph(args.graph\_type, args.n\_nodes, seed}\OperatorTok{=}\NormalTok{args.seed)}
\NormalTok{    H\_C }\OperatorTok{=}\NormalTok{ build\_cost\_hamiltonian(graph)}
    \CommentTok{\# Brute{-}force C\_max for N \textless{}= 20}
    \ControlFlowTok{if}\NormalTok{ args.n\_nodes }\OperatorTok{\textless{}=} \DecValTok{20}\NormalTok{:}
\NormalTok{        c\_max }\OperatorTok{=}\NormalTok{ brute\_force\_max\_cut(graph)}
\NormalTok{    ...}

    \CommentTok{\# 2. Device selection}
    \ControlFlowTok{if}\NormalTok{ args.hardware:}
\NormalTok{        dev }\OperatorTok{=}\NormalTok{ create\_ibm\_device(args.n\_nodes, args.backend)}
    \ControlFlowTok{elif}\NormalTok{ args.noise:}
\NormalTok{        dev }\OperatorTok{=}\NormalTok{ create\_noisy\_device(args.n\_nodes, ...)}
    \ControlFlowTok{else}\NormalTok{:}
\NormalTok{        dev }\OperatorTok{=}\NormalTok{ qml.device(}\StringTok{"default.qubit"}\NormalTok{, wires}\OperatorTok{=}\NormalTok{args.n\_nodes)}

\NormalTok{    qnode }\OperatorTok{=}\NormalTok{ build\_qaoa\_circuit(H\_C, args.n\_nodes, args.p, dev)}

    \CommentTok{\# 3. Landscape sampling (p=1 with rendering)}
    \ControlFlowTok{if}\NormalTok{ args.render }\KeywordTok{and}\NormalTok{ args.p }\OperatorTok{==} \DecValTok{1}\NormalTok{:}
\NormalTok{        landscape }\OperatorTok{=}\NormalTok{ sample\_landscape\_parallel(}
\NormalTok{            qnode, args.p, (}\DecValTok{0}\NormalTok{, np.pi), (}\DecValTok{0}\NormalTok{, np.pi}\OperatorTok{/}\DecValTok{2}\NormalTok{),}
\NormalTok{            args.grid\_size, args.grid\_size, c\_max)}

    \CommentTok{\# 4. Classical optimisation}
\NormalTok{    result }\OperatorTok{=}\NormalTok{ optimize\_qaoa(qnode, args.p, method}\OperatorTok{=}\NormalTok{args.optimizer, ...)}

    \CommentTok{\# 5. ZNE error mitigation}
    \ControlFlowTok{if}\NormalTok{ args.zne }\KeywordTok{and}\NormalTok{ (args.noise }\KeywordTok{or}\NormalTok{ args.hardware):}
\NormalTok{        E0, details }\OperatorTok{=}\NormalTok{ zne\_extrapolate(qnode, result.best\_params,}
\NormalTok{                                       scale\_factors}\OperatorTok{=}\NormalTok{[}\DecValTok{1}\NormalTok{, }\DecValTok{3}\NormalTok{, }\DecValTok{5}\NormalTok{])}

    \CommentTok{\# 6. Rendering}
    \ControlFlowTok{if}\NormalTok{ args.render:}
\NormalTok{        window }\OperatorTok{=}\NormalTok{ init\_opengl\_context(}\DecValTok{1280}\NormalTok{, }\DecValTok{720}\NormalTok{, }\StringTok{"QAOA Landscape"}\NormalTok{)}
\NormalTok{        vao, vbo, ebo, n\_idx }\OperatorTok{=}\NormalTok{ create\_surface\_vbo(}
\NormalTok{            landscape.gamma\_vals, landscape.beta\_vals, landscape.energies)}
        \CommentTok{\# Main render loop}
        \ControlFlowTok{while} \KeywordTok{not}\NormalTok{ glfw.window\_should\_close(window):}
\NormalTok{            glClear(GL\_COLOR\_BUFFER\_BIT }\OperatorTok{|}\NormalTok{ GL\_DEPTH\_BUFFER\_BIT)}
            \CommentTok{\# Draw surface, trajectory, side panel}
\NormalTok{            ...}
\NormalTok{            glfw.swap\_buffers(window)}
\NormalTok{            glfw.poll\_events()}
\NormalTok{        glfw.terminate()}
\end{Highlighting}
\end{Shaded}

\emph{(Full implementation: \texttt{report1\_scripts/main.py})}

\begin{center}\rule{0.5\linewidth}{0.5pt}\end{center}

\section{Part 4 --- OpenGL Visualization
Design}\label{part-4-opengl-visualization-design}

\subsection{\texorpdfstring{1. Full Rendering Pipeline for the 3D
\((\beta, \gamma, \langle H_C \rangle)\)
Surface}{1. Full Rendering Pipeline for the 3D (\textbackslash beta, \textbackslash gamma, \textbackslash langle H\_C \textbackslash rangle) Surface}}\label{full-rendering-pipeline-for-the-3d-beta-gamma-langle-h_c-rangle-surface}

\subsubsection{Vertex Buffer Layout}\label{vertex-buffer-layout}

Each vertex has 7 float32 values (28 bytes per vertex):

{\def\LTcaptype{none} % do not increment counter
\begin{longtable}[]{@{}
  >{\centering\arraybackslash}p{(\linewidth - 6\tabcolsep) * \real{0.2500}}
  >{\centering\arraybackslash}p{(\linewidth - 6\tabcolsep) * \real{0.2500}}
  >{\centering\arraybackslash}p{(\linewidth - 6\tabcolsep) * \real{0.2500}}
  >{\raggedright\arraybackslash}p{(\linewidth - 6\tabcolsep) * \real{0.2500}}@{}}
\toprule\noalign{}
\begin{minipage}[b]{\linewidth}\centering
Offset (bytes)
\end{minipage} & \begin{minipage}[b]{\linewidth}\centering
Attribute
\end{minipage} & \begin{minipage}[b]{\linewidth}\centering
Size
\end{minipage} & \begin{minipage}[b]{\linewidth}\raggedright
Description
\end{minipage} \\
\midrule\noalign{}
\endhead
\bottomrule\noalign{}
\endlastfoot
0 & Position & 3 floats &
\((x, y, z) = (\gamma, \beta, \langle H_C \rangle)\) \\
12 & Normal & 3 floats & Surface normal \((n_x, n_y, n_z)\) for Phong
lighting \\
24 & Energy & 1 float & Normalised energy \(\in [0, 1]\) for colour
mapping \\
\end{longtable}
}

The surface mesh is a regular grid of \((N_\gamma \times N_\beta)\)
vertices connected by \((N_\gamma - 1)(N_\beta - 1) \times 2\) triangles
(two per grid cell). The index buffer stores triangle indices as uint32.

\subsubsection{GLSL Shaders}\label{glsl-shaders}

\textbf{Vertex Shader (\texttt{surface.vert}):}

The vertex shader receives three per-vertex attributes at layout
locations 0--2: position \((\gamma, \beta, \langle H_C\rangle)\), the
surface normal, and the normalised energy scalar. It transforms the
position into world space via the model matrix, computes the correct
world-space normal using the inverse-transpose of the model matrix (to
handle non-uniform scaling), passes the normalised energy through as a
varying, and outputs clip-space position via the MVP chain.

\emph{(Full source: \texttt{report1\_scripts/shaders/surface.vert})}

\textbf{Fragment Shader (\texttt{surface.frag}):}

The fragment shader implements Phong illumination with an energy-based
colour map. The normalised energy \(e \in [0,1]\) is mapped to a hue via
\(h = (1 - e) \times 0.66\), producing a blue-to-red gradient (blue =
low energy, red = high energy) through an HSV→RGB conversion. Phong
shading combines three terms:

\begin{itemize}
\tightlist
\item
  \textbf{Ambient} (\(0.2 \times\) base colour) --- constant fill light.
\item
  \textbf{Diffuse} (\(\max(\hat{n} \cdot \hat{l},\, 0) \times\) base
  colour) --- Lambertian term.
\item
  \textbf{Specular}
  (\(0.5 \times [\max(\hat{v} \cdot \hat{r},\, 0)]^{32}\)) --- shininess
  highlights.
\end{itemize}

The output alpha is \(0.9\) (slight transparency) to allow seeing the
trajectory through the surface edges.

\emph{(Full source: \texttt{report1\_scripts/shaders/surface.frag})}

\subsubsection{Camera Controls (Arcball
Rotation)}\label{camera-controls-arcball-rotation}

The \texttt{ArcballCamera} class (Module 8) handles: -
\textbf{Left-click drag:} Rotate the view around the scene centre. The
yaw/pitch angles update proportionally to mouse displacement. -
\textbf{Scroll wheel:} Zoom in/out by adjusting the camera-to-target
distance. - \textbf{Middle-click drag:} Pan the target point in the view
plane.

GLFW callbacks wire mouse events to the camera. A cursor-position
callback computes the delta from the last position and calls
\texttt{camera.on\_mouse\_drag(dx,\ dy)} when the left button is held. A
scroll callback forwards the vertical offset to
\texttt{camera.on\_scroll(yoffset)}. Both callbacks are registered via
\texttt{glfw.set\_cursor\_pos\_callback} and
\texttt{glfw.set\_scroll\_callback}.

\emph{(Full implementation: \texttt{report1\_scripts/renderer.py})}

\subsubsection{Live Surface Update During
Optimization}\label{live-surface-update-during-optimization}

As the optimizer progresses, the landscape surface can be updated in two
ways:

\begin{enumerate}
\def\labelenumi{\arabic{enumi}.}
\item
  \textbf{Full recompute (for changing landscape slice in \(p > 1\)):}
  When the fixed parameters change, resample the 2D landscape and call
  \texttt{update\_surface\_vbo} with new energy data. This uses
  \texttt{glBufferSubData} for an in-place update. The recompute is
  throttled (e.g., every 10 optimiser steps) to avoid bottlenecking the
  render loop.
\item
  \textbf{Streaming VBO update (double buffering):} Two VBOs are
  allocated. While one is bound for rendering, the other is updated on a
  background thread. After the update completes, the buffer IDs are
  swapped. This prevents stalls in the render loop during data upload.
\end{enumerate}

\subsection{2. Optimizer Trajectory
Overlay}\label{optimizer-trajectory-overlay}

The optimizer trajectory is rendered as a 3D polyline lying on (or
slightly above) the energy surface. Each vertex is
\((\gamma_k, \beta_k, \langle H_C \rangle_k + \epsilon)\) where
\(\epsilon \approx 0.01\) prevents z-fighting.

\textbf{Trajectory shaders:}

The trajectory vertex shader (\texttt{trajectory.vert}) takes position
and a progress scalar \(t \in [0,1]\) as inputs, applies the combined
model--view--projection matrix, and forwards \(t\) to the fragment
stage. The fragment shader (\texttt{trajectory.frag}) interpolates
colour from blue \((0.0, 0.3, 1.0)\) at \(t = 0\) to red
\((1.0, 0.1, 0.0)\) at \(t = 1\) using GLSL's \texttt{mix} function,
with full opacity.

\emph{(Full source: \texttt{report1\_scripts/shaders/trajectory.vert}
and \texttt{trajectory.frag})}

The trajectory VBO is rebuilt every frame (or every N iterations) from
the \texttt{trajectory\_log} list using \texttt{GL\_STREAM\_DRAW}.

\subsection{3. Side Panel}\label{side-panel}

Using Dear ImGui (via \texttt{imgui{[}glfw{]}} Python bindings) for the
info panel. Each frame, a fixed-size window is positioned in the
top-left corner and populated with live text fields: circuit depth
\(p\), current \(\langle \hat{H}_C \rangle\), approximation ratio,
gradient norm, iteration count, optimiser name, and graph metadata. The
ImGui draw data is rendered to the OpenGL framebuffer after the 3D
scene.

\emph{(Full implementation: \texttt{report1\_scripts/renderer.py})}

\begin{center}\rule{0.5\linewidth}{0.5pt}\end{center}

\section{Part 5 --- Hardware Execution
Plan}\label{part-5-hardware-execution-plan}

\subsection{1. IBM Quantum Submission via
PennyLane}\label{ibm-quantum-submission-via-pennylane}

\subsubsection{1.1 Transpilation for a Specific
Backend}\label{transpilation-for-a-specific-backend}

PennyLane's \texttt{qiskit.ibmq} device handles transpilation
internally. When the QNode is executed, PennyLane:

\begin{enumerate}
\def\labelenumi{\arabic{enumi}.}
\tightlist
\item
  Converts PennyLane operations to Qiskit gates.
\item
  Calls \texttt{qiskit.transpile()} targeting the backend's native gate
  set (typically \(\{CX, \text{ID}, R_Z, \sqrt{X}, X\}\) for IBM
  Eagle/Heron processors).
\item
  Maps virtual qubits to physical qubits based on the coupling map.
\item
  Applies routing passes (SABRE, stochastic swap) to handle non-adjacent
  CNOT requirements.
\end{enumerate}

To control transpilation, the PennyLane-Qiskit device accepts a
\texttt{transpile\_options} dictionary specifying
\texttt{optimization\_level} (0--3, where 3 enables the most aggressive
gate cancellation and resynthesis passes) and \texttt{initial\_layout}
(the list of physical qubit indices returned by
\texttt{select\_optimal\_qubits}). These are forwarded directly to
\texttt{qiskit.transpile()}.

\emph{(Full implementation:
\texttt{report1\_scripts/hardware\_runner.py})}

\subsubsection{1.2 Qubit Selection via Coupling Map
Analysis}\label{qubit-selection-via-coupling-map-analysis}

Use \texttt{select\_optimal\_qubits()} (Module 6). The key heuristic:

\begin{enumerate}
\def\labelenumi{\arabic{enumi}.}
\tightlist
\item
  Load calibration data for all qubits (T1, T2) and all CNOT/ECR gates
  (error rates).
\item
  Build a weighted graph over physical qubits.
\item
  Find the connected subgraph of size \(N\) that maximises
  \(\sum_{\text{edges}} 1/\epsilon_{CX} + \alpha \sum_{\text{nodes}} T_1 T_2\).
\item
  Use this subgraph as the \texttt{initial\_layout} in transpilation.
\end{enumerate}

For the QAOA circuit on \(C_4\) (4 qubits, ring topology), select 4
physical qubits that form a ring in the coupling map with minimal CNOT
error.

\subsubsection{\texorpdfstring{1.3 Shots: \texttt{qml.sample} vs
\texttt{qml.expval}}{1.3 Shots: qml.sample vs qml.expval}}\label{shots-qml.sample-vs-qml.expval}

\textbf{\texttt{qml.expval(H\_C)} on hardware:} PennyLane decomposes
\(H_C = \sum_k c_k P_k\) into Pauli terms, groups commuting terms, and
measures each group separately. For MaxCut on \(|E|\) edges, \(H_C\) has
\(|E|\) ZZ terms --- all mutually commuting (all diagonal), so they can
be measured in a single circuit run. This is optimal.

\textbf{\texttt{qml.sample()} on hardware:} Returns raw bitstrings,
shape \texttt{(shots,\ n\_qubits)}. You compute \(\langle H_C \rangle\)
classically:

\[\langle H_C \rangle \approx \frac{1}{S}\sum_{s=1}^{S} C(\mathbf{z}^{(s)})\]

where \(C(\mathbf{z})\) is the MaxCut cost evaluated on bitstring
\(\mathbf{z}\). This avoids PennyLane's Pauli grouping overhead and
gives access to the full distribution for post-processing.

\textbf{Shot budget:} For standard deviation \(\sigma\) in the energy
estimate:

\[\sigma = \frac{\Delta E}{\sqrt{S}}\]

where \(\Delta E = E_{\max} - E_{\min}\) is the spectral range of
\(H_C\). To achieve precision \(\epsilon\), need
\(S \geq (\Delta E / \epsilon)^2\) shots. For MaxCut on 4-10 qubits,
\(\Delta E \sim 10\), so for \(\epsilon = 0.1\): \(S \geq 10000\) shots.

\subsubsection{1.4 Readout Error Handling}\label{readout-error-handling}

Readout (measurement) errors are the dominant noise source on IBM
hardware (typically 1-5\% per qubit). Strategies:

\begin{enumerate}
\def\labelenumi{\arabic{enumi}.}
\item
  \textbf{Readout error mitigation:} Prepare all \(2^N\) computational
  basis states, measure, and build a calibration matrix \(M\) where
  \(M_{ij} = P(\text{measure } i | \text{prepared } j)\). Then invert:
  \(\mathbf{p}_{\text{corrected}} = M^{-1} \mathbf{p}_{\text{raw}}\).
\item
  \textbf{Twirled readout error mitigation (TREX):} Randomly flip
  measurement outcomes using Pauli-\(X\) twirling to convert general
  readout error into a symmetric channel, then correct analytically.
\item
  \textbf{PennyLane integration:} Use
  \texttt{qml.transforms.mitigate\_with\_zne} or custom post-processing
  on the raw counts.
\end{enumerate}

\subsection{2. Error Mitigation
Strategy}\label{error-mitigation-strategy}

\subsubsection{2.1 ZNE with Linear + Richardson
Extrapolation}\label{zne-with-linear-richardson-extrapolation}

\textbf{Global unitary folding:} To amplify noise by factor \(\lambda\),
replace circuit \(U\) with \(U(U^\dagger U)^{(\lambda-1)/2}\):

\begin{itemize}
\tightlist
\item
  \(\lambda = 1\): original circuit (depth \(d\))
\item
  \(\lambda = 3\): circuit followed by \(U^\dagger U\) (depth \(3d\))
\item
  \(\lambda = 5\): circuit followed by \(U^\dagger U U^\dagger U\)
  (depth \(5d\))
\end{itemize}

For each \(\lambda\), measure \(\langle H_C \rangle_\lambda\).

\textbf{Linear extrapolation:} Fit \(E(\lambda) = a + b\lambda\) and
extrapolate to \(\lambda = 0\): \(E_0 = a = E(\lambda=1) - b\). Assumes
noise effect is linear in noise strength.

\textbf{Richardson extrapolation:} Fit
\(E(\lambda) = \sum_{k=0}^{M-1} a_k \lambda^k\) through \(M\) data
points and evaluate at \(\lambda = 0\). More accurate but more sensitive
to statistical noise.

\textbf{ZNE for QAOA MaxCut on \(C_4\) at \(p = 1\):}

With the circuit depth being \(|E| + N = 4 + 4 = 8\) gates (4
ZZ-rotations + 4 RX gates): - \(\lambda = 1\): 8 gates -
\(\lambda = 3\): 24 gates - \(\lambda = 5\): 40 gates

Use 3 scale factors with Richardson extrapolation by calling
\texttt{zne\_extrapolate(hardware\_qnode,\ optimal\_params,\ scale\_factors={[}1,\ 3,\ 5{]},\ extrapolation="richardson")}.

\emph{(Full implementation:
\texttt{report1\_scripts/error\_mitigation.py})}

\subsubsection{2.2 Pauli Twirling as
Precursor}\label{pauli-twirling-as-precursor}

Pauli twirling converts general coherent errors on 2-qubit gates into
stochastic Pauli channels (incoherent noise), which ZNE handles more
effectively. For each CNOT (or ECR) gate in the transpiled circuit:

\begin{enumerate}
\def\labelenumi{\arabic{enumi}.}
\tightlist
\item
  Sample random Pauli pair \((P_a, P_b)\) from the twirl group.
\item
  Replace \(\text{CNOT}\) with \(P_b \cdot \text{CNOT} \cdot P_a\).
\item
  Average over many random twirls.
\end{enumerate}

This should be applied before ZNE. PennyLane does not have built-in
Pauli twirling, so it must be implemented via tape transforms.

\subsubsection{2.3 Realistic Fidelity
Expectations}\label{realistic-fidelity-expectations}

{\def\LTcaptype{none} % do not increment counter
\begin{longtable}[]{@{}
  >{\raggedright\arraybackslash}p{(\linewidth - 4\tabcolsep) * \real{0.3333}}
  >{\centering\arraybackslash}p{(\linewidth - 4\tabcolsep) * \real{0.3333}}
  >{\centering\arraybackslash}p{(\linewidth - 4\tabcolsep) * \real{0.3333}}@{}}
\toprule\noalign{}
\begin{minipage}[b]{\linewidth}\raggedright
Configuration
\end{minipage} & \begin{minipage}[b]{\linewidth}\centering
Expected \(\langle H_C \rangle\) error
\end{minipage} & \begin{minipage}[b]{\linewidth}\centering
Approx. ratio degradation
\end{minipage} \\
\midrule\noalign{}
\endhead
\bottomrule\noalign{}
\endlastfoot
\(p=1\), 4 qubits, noiseless & 0 & 0 \\
\(p=1\), 4 qubits, ibm\_brisbane (raw) & \textasciitilde0.2-0.5 &
\textasciitilde5-15\% \\
\(p=1\), 4 qubits, ibm\_brisbane + ZNE & \textasciitilde0.05-0.15 &
\textasciitilde2-5\% \\
\(p=1\), 10 qubits, ibm\_brisbane (raw) & \textasciitilde0.5-1.5 &
\textasciitilde10-30\% \\
\(p=1\), 10 qubits, ibm\_brisbane + ZNE & \textasciitilde0.2-0.5 &
\textasciitilde5-15\% \\
\(p=2\), 10 qubits, ibm\_brisbane (raw) & \textasciitilde1.0-3.0 &
\textasciitilde20-50\% \\
\(p=2\), 10 qubits, ibm\_brisbane + ZNE & \textasciitilde0.3-1.0 &
\textasciitilde10-25\% \\
\end{longtable}
}

These are order-of-magnitude estimates based on typical 2024-2025 IBM
Eagle processor performance with \(\sim 1\%\) two-qubit gate error rates
and \(\sim 2\%\) readout error rates.

\subsection{3. Benchmarking Checklist}\label{benchmarking-checklist}

\textbf{Metrics to record for each experiment:}

{\def\LTcaptype{none} % do not increment counter
\begin{longtable}[]{@{}
  >{\raggedright\arraybackslash}p{(\linewidth - 8\tabcolsep) * \real{0.2000}}
  >{\centering\arraybackslash}p{(\linewidth - 8\tabcolsep) * \real{0.2000}}
  >{\centering\arraybackslash}p{(\linewidth - 8\tabcolsep) * \real{0.2000}}
  >{\centering\arraybackslash}p{(\linewidth - 8\tabcolsep) * \real{0.2000}}
  >{\centering\arraybackslash}p{(\linewidth - 8\tabcolsep) * \real{0.2000}}@{}}
\toprule\noalign{}
\begin{minipage}[b]{\linewidth}\raggedright
Metric
\end{minipage} & \begin{minipage}[b]{\linewidth}\centering
Noiseless Sim
\end{minipage} & \begin{minipage}[b]{\linewidth}\centering
Noisy Sim
\end{minipage} & \begin{minipage}[b]{\linewidth}\centering
Hardware
\end{minipage} & \begin{minipage}[b]{\linewidth}\centering
Hardware+ZNE
\end{minipage} \\
\midrule\noalign{}
\endhead
\bottomrule\noalign{}
\endlastfoot
\(\langle H_C \rangle\) & ✓ & ✓ & ✓ & ✓ \\
Approximation ratio \(r = \langle H_C \rangle / C_{\max}\) & ✓ & ✓ & ✓ &
✓ \\
Optimal \((\gamma^*, \beta^*)\) & ✓ & ✓ & ✓ & ✓ \\
Optimizer iterations to convergence & ✓ & ✓ & ✓ & ✓ \\
Total circuit evaluations & ✓ & ✓ & ✓ & ✓ \\
Wall-clock time & ✓ & ✓ & ✓ & ✓ \\
Bitstring distribution (top-10 samples) & ✓ & ✓ & ✓ & ✓ \\
Probability of sampling optimal cut & ✓ & ✓ & ✓ & ✓ \\
\end{longtable}
}

\textbf{Plots:} 1. \(r\) vs \(p\) (circuit depth): noiseless, noisy,
hardware, hardware+ZNE on same axes. 2. Energy landscape surface
comparison: noiseless vs noisy (as 3D surfaces or heatmaps side by
side). 3. Optimizer convergence curves: energy vs iteration for each
configuration. 4. Histogram of sampled bitstring costs: noiseless vs
hardware.

\begin{center}\rule{0.5\linewidth}{0.5pt}\end{center}

\section{Part 6 --- Research Contribution
Angle}\label{part-6-research-contribution-angle}

\subsection{1. Three Directions Toward a Publishable
Paper}\label{three-directions-toward-a-publishable-paper}

\textbf{Direction 1: QAOA Landscape Phase Transitions in Frustrated
Graphs}

Study the QAOA energy landscape topology as a function of graph
frustration index. For anti-ferromagnetic Ising models on non-bipartite
graphs (triangular lattices, Petersen graph, random \(d\)-regular), map
the landscape at \(p = 1, 2, 3\) and quantify: - Number of local minima
(via landscape sampling + basin detection). - Hessian spectrum at
critical points (positive definite, negative definite, saddle index). -
Correlation between landscape complexity and classical hardness
(brute-force time, simulated annealing residual energy).

This extends the 2025 Boulebnane-Montanaro result by moving beyond
regular graphs with uniform fields to arbitrary frustrated topologies.

\textbf{Direction 2: Noise-Induced Landscape Smoothing and
Concentration}

Systematically compare the \((\beta, \gamma)\) landscape under
noiseless, depolarising, and thermal relaxation noise at varying
strengths. The hypothesis: moderate noise smooths away spurious local
minima, potentially improving optimizer convergence (stochastic
regularisation). Measure: - Gradient norm distribution across the
landscape as a function of noise rate. - Number of local minima as noise
increases (using persistent homology / topological data analysis on the
landscape). - Optimal parameter stability: how much do
\((\gamma^*, \beta^*)\) shift under noise?

This connects to the barren plateau literature and provides practical
guidance on when noise helps vs hurts variational optimization.

\textbf{Direction 3: Adaptive QAOA Depth Selection via Landscape
Diagnostics}

Propose a protocol that examines the \(p\)-level landscape structure to
decide whether increasing \(p\) is worthwhile: - Compute the landscape
Hessian at the \(p\)-level optimum. - If the smallest Hessian eigenvalue
is ``sufficiently negative'' (surface is steeply curved), the current
\(p\) is under-parameterised and increasing \(p\) will help. - If the
Hessian is nearly flat (plateau), further depth is wasteful. - Validate
on families of graphs with known \(p^*\) (the smallest \(p\) for which
QAOA is provably optimal).

\subsection{2. Most Relevant 2024--2025
Papers}\label{most-relevant-20242025-papers}

\begin{enumerate}
\def\labelenumi{\arabic{enumi}.}
\item
  \textbf{Boulebnane \& Montanaro (2025).} ``Solving optimization
  problems with local light cones on fixedangle QAOA.'' \emph{Quantum}.
  --- Phase transition in QAOA performance with longitudinal field.
\item
  \textbf{Kremenetski, Hogg, \& Hadfield (2025).} ``Quantum Alternating
  Operator Ansatz (QAOA+): landscape structure with constraints.''
  \emph{PRX Quantum}. --- Extended QAOA landscape analysis beyond
  MaxCut.
\item
  \textbf{Shaydulin et al.~(2024).} ``Evidence of scaling advantage for
  the quantum approximate optimization algorithm on a classically
  intractable problem.'' \emph{Science Advances}. --- First empirical
  evidence of QAOA advantage at scale.
\item
  \textbf{Wurtz \& Love (2024).} ``Counterdiabatically optimized QAOA.''
  \emph{Physical Review A}. --- Improved initialisation strategies via
  counterdiabatic terms.
\item
  \textbf{Cerezo et al.~(2024).} ``Does provable absence of barren
  plateaus imply classical simulability?'' \emph{Physical Review X}. ---
  Fundamental limits on trainability and simulability.
\item
  \textbf{Pelofske et al.~(2024).} ``Scaling whole-chip QAOA for
  higher-order ising spin glass models on heavy-hex graphs.'' \emph{npj
  Quantum Information}. --- Large-scale QAOA benchmarking on IBM
  hardware.
\item
  \textbf{Zhou, Wang, Choi et al.~(2020).} ``Quantum Approximate
  Optimization Algorithm: Performance, Mechanism, and Implementation on
  Near-Term Devices.'' \emph{Physical Review X}. --- INTERP and FOURIER
  strategies (seminal reference).
\end{enumerate}

\subsection{3. Interesting Graph Families and Problem
Types}\label{interesting-graph-families-and-problem-types}

{\def\LTcaptype{none} % do not increment counter
\begin{longtable}[]{@{}
  >{\raggedright\arraybackslash}p{(\linewidth - 4\tabcolsep) * \real{0.3333}}
  >{\raggedright\arraybackslash}p{(\linewidth - 4\tabcolsep) * \real{0.3333}}
  >{\raggedright\arraybackslash}p{(\linewidth - 4\tabcolsep) * \real{0.3333}}@{}}
\toprule\noalign{}
\begin{minipage}[b]{\linewidth}\raggedright
Graph Family
\end{minipage} & \begin{minipage}[b]{\linewidth}\raggedright
Why Interesting
\end{minipage} & \begin{minipage}[b]{\linewidth}\raggedright
Landscape Feature
\end{minipage} \\
\midrule\noalign{}
\endhead
\bottomrule\noalign{}
\endlastfoot
\textbf{Random 3-regular} & Phase transition at clause density for
SAT-like problems & Complex landscape with many near-degenerate
minima \\
\textbf{Petersen graph} & Highly symmetric, non-planar, frustrated &
Discrete symmetry in landscape, symmetry-breaking effects \\
\textbf{Triangular lattice} & Geometrically frustrated for AF Ising &
Exponentially degenerate ground states, flat landscape regions \\
\textbf{Erdős-Rényi \(G(n, p_c)\)} & Percolation transition at
\(p_c = 1/n\) & Landscape topology changes at the percolation
threshold \\
\textbf{Complete bipartite \(K_{n,n}\)} & MaxCut is trivial but
landscape is instructive & Single clean global minimum --- contrast with
random graphs \\
\textbf{Weighted Sherrington-Kirkpatrick} & Complete graph with Gaussian
random weights & Full-RSB spin glass: hardest possible QAOA landscape \\
\end{longtable}
}

The most visually dramatic landscapes come from frustrated systems
(triangular, 3-regular non-bipartite, SK model) where near-degenerate
minima create complex topography.

\end{document}
